\section{CM fields with character and Selmer bounds}
\label{sec:family}

Fix a primitive Dirichlet character \(\psi\) of conductor \(M\) and a
real exponent parameter \(0<\varepsilon<1/52\). The zeta case takes
\(\psi=1\) of conductor \(1\). For every nonprincipal Dirichlet
character \(\chi\), \(L(s,\chi)\) is holomorphic at \(s=1\) and
\(L(1,\chi)\ne0\) \cite[\S25.15(i) and (25.15.9)]{dlmf}. For
primitive nonprincipal \(\chi\) define
\begin{equation}
\label{eq:centred-prime-sums}
  G_\chi(y)=\sum_{n\ge1}\Lambda(n)\chi(n)(y/n-1)_+
                  +y\frac{L'(1,\chi)}{L(1,\chi)}
  \qquad(y>0).
\end{equation}
For \(\chi=1\), use the zeta sum \(G_1\) of \eqref{eq:principal-ramp}.

At each real scale \(x\), the local division estimates require bounds
on Selmer groups over a CM field \(E\). The Fourier estimate requires
prime-sum bounds for characters of \(\operatorname{Gal}(E/\mathbf Q)\).
The two bounds hold for the same choice of \(E\) at each point
of the intervals used to construct the family. The classical
lattice and Fourier comparison will evaluate that family at
\hyperlink{proof.classical-window-test-scale}{one scale in each window}.

\subsection{The auxiliary primes and character estimates}

At an auxiliary scale \(Z\), put \(w=Z^{\varepsilon/5}\). The primes
\(q,d\) are distinct and coprime to \(M\), with
\[
\begin{gathered}
 q\equiv3\pmod 8,\qquad d\equiv7\pmod {16},\qquad
 q\mid d+1,\qquad p\mid d+1\quad(3\le p\le w,\ p\text{ prime}),\\
 \tfrac12w<\log q\le w,\qquad
 \tfrac12w(\log Z)^2<\log d\le w(\log Z)^2 .
\end{gathered}
\]
Write \(\eta_q=(\frac{\cdot}{q})\) and \(\eta_d=(\frac{\cdot}{d})\) for
the quadratic Dirichlet characters. For \(h\ge0\) let
\[
 E(h,q,d)=\mathbf Q(\zeta_{2^{h+2}})^+\,\mathbf Q(\psi)\,
              \mathbf Q(\sqrt{-q},\sqrt{-d}),
\]
where \(\mathbf Q(\psi)\) is the subfield of \(\mathbf Q(\zeta_M)\)
fixed by the kernel of \(\psi\). Since \(E(h,q,d)\) is abelian and
contains \(\mathbf Q(\sqrt{-q})\), it is CM.

For an odd or \(4\)-divisible modulus \(N\), the cyclotomic Artin map
identifies \((\mathbf Z/N\mathbf Z)^\times\) with
\(\operatorname{Gal}(\mathbf Q(\zeta_N)/\mathbf Q)\) by
\(a\mapsto(\zeta_N\mapsto\zeta_N^a)\). At any prime \(\ell\nmid N\)
the residue class of \(\ell\) gives the arithmetic Frobenius.
For a prime \(p\equiv3\bmod4\), the unique quadratic subfield of
\(\mathbf Q(\zeta_p)\) is \(\mathbf Q(\sqrt{-p})\), and the
restriction of the automorphism indexed by \(a\bmod p\) to it acts
on \(\sqrt{-p}\) by the sign \((\frac ap)\).
\cite[Ch.~V, Examples~3.1--3.2 and 3.11]{f1milne-cft}
Applying this to \(q,d\) gives the quadratic factors of conductors
\(q,d\). These are coprime to \(\operatorname{lcm}(2^{h+2},M)\) and to
each other, so the Chinese remainder theorem gives
\(\Gamma=\operatorname{Gal}(E(h,q,d)/\mathbf Q)\simeq
(\mathbf Z/2\mathbf Z)^2\times H\), with \(H\) the group of the
compositum of the first two factors.
For \(a\) coprime to \(2Mqd\), let \(h_{\rm nat}(a)\in H\) be the
restriction of the cyclotomic automorphism \(\zeta\mapsto\zeta^a\)
to the compositum. Put \(h_-=h_{\rm nat}(-1)\), so \(2h_-=0\).
Write \(\eta_d(a)=(-1)^{e_d(a)}\) and \(\eta_q(a)=(-1)^{e_q(a)}\).
The release's
\releasecite{Formalism/SharpRestriction/RamificationFrobenius/TotientGaloisClasses/InertiaGaloisSymmetry/AnyFrame.lean}{235}{Artin map}
orders the signs as \(q,d\) and shifts the last component to
\(h_{\rm nat}(a)+e_q(a)h_-\). The
\releasecite{Development/AutomorphicArithmetic/PrimitiveCosets/RecH.lean}{10}{Fourier coordinates}
interchange the two sign positions, giving
\begin{equation}
\label{eq:artin-coordinates}
 \operatorname{Art}(a)
 =\bigl(e_d(a),e_q(a),h_{\rm nat}(a)+e_q(a)h_-\bigr).
\end{equation}
Both sign bits of \(-1\) are \(1\), so complex conjugation has
coordinates \(c_\infty=(1,1,0)\). For \(\psi=1\) the compositum is
totally real and \(h_-=0\).

For odd primes \(\ell\nmid d\), splitting in \(\mathbf Q(\sqrt{-d})\)
is equivalent to \(-d\) being a square modulo \(\ell\). For \(\ell=2\)
it is equivalent to \(-d\equiv1\pmod8\), since \(-d\) is an odd
quadratic discriminant. Thus \(2\) and the specified odd primes
dividing \(d+1\) split. These include \(q\) and, once \(w\ge4M\),
every odd prime dividing \(M\).

\paragraph{The classical host and its Artin data.}
For \(\psi=1\), put \(Q_S=2^{h+2}\) and \(Q=Q_Sqd\). The field
\(E=E(h,q,d)\) is the fixed field in \(\mathbf Q(\zeta_Q)\) of the
units \(a\) satisfying \(a\equiv\pm1\bmod Q_S\) and
\(\eta_q(a)=\eta_d(a)=1\). Indeed, these are precisely the
automorphisms fixing each of its three factors. This also identifies
any field of the prescribed type in the Selmer assertion with \(E\).
In an ambient cyclotomic field of a multiple of \(Q\), its defining
subgroup is the inverse image of the same subgroup modulo \(Q\).
The group \(H=(\mathbf Z/Q_S\mathbf Z)^\times/\{\pm1\}\) has order
\(2^h\) and is cyclic. The class of \(5\) generates it, since
\(v_2(5^{2^j}-1)=j+2\) by repeated squaring and no power of \(5\)
is \(-1\bmod 4\). Here \(H=H_0\). The fixed field \(F=E^{\langle
c_\infty\rangle}\) is totally real and \([F:\mathbf Q]=2|H|\).

For \(p\mid Q\), the \(p\)-power cyclotomic factor is totally ramified
locally, by the Eisenstein polynomial of \(\zeta_{p^a}-1\) when
\(p^a\Vert Q\). The roots of unity of order prime to \(p\) lie in the
maximal unramified extension, by lifting the roots in the residue
field. Arithmetic Frobenius raises these roots to their \(p\)-th
powers. Thus the inertia subgroup is the image of the \(p\)-factor
of the group of units in the Chinese-remainder decomposition. In
the coordinates of \eqref{eq:artin-coordinates} this gives
\begin{equation}
\label{eq:classical-host-inertia}
 I_2=\{(0,0,s):s\in H\},\qquad
 I_q=\{0,(0,1,0)\},\qquad I_d=\{0,(1,0,0)\},
 \qquad I_p=\{0\}\quad(p\nmid Q).
\end{equation}
None of these contains \(c_\infty=(1,1,0)\), so \(E/F\) is
unramified at all finite places. Put \(c(e_d,e_q,s)=(-1)^{e_d}\)
and \(P=\ker c\). Both \(2\) and \(q\) split in
\(\mathbf Q(\sqrt{-d})\). Therefore every Frobenius lift at \(2\)
or \(q\), as well as its inertia subgroup, lies in \(P\).
For every lift \(F_d\) and every integer \(j\ge1\), the coset
\(jF_d+I_d\) has one element outside \(P\), since \(c\) is
nontrivial on \(I_d\). At \(p\nmid Q\) the Frobenius is
\(\operatorname{Art}(p)\). These assertions give the local inputs
to \eqref{eq:prime-power-ramp} and its Fourier expansion. The character
\(c\) induces the primitive character \(\eta_d\).

\teamanchor{proof.classical-host-roots-of-unity}%
For large scales \(q,d>3\), the roots of unity in \(E\) are
\(\{1,-1\}\). If an odd prime \(\ell\) divides the order of a
root of unity in \(E\), then \(\mathbf Q(\zeta_\ell)\subset E\).
Its ramification at \(\ell\) forces \(\ell=q\) or \(d\), by
\eqref{eq:classical-host-inertia}. Its inertia order \(\ell-1\)
would then divide \(2\), which is impossible. The unit class equal
to \(-1\bmod Q_S\) and \(1\bmod qd\) fixes \(E\) in
\(\mathbf Q(\zeta_Q)\) and acts nontrivially on \(\zeta_4\).
Thus \(E\) contains no root of unity of order divisible by \(4\).

For \(h\ge1\) there is just one nontrivial real character \(\nu\)
of \(H\). On odd integers it is \((2/\cdot)\), since the class of
\(5\) has value \(-1\) and the map through units modulo \(8\) is even.
Its primitive conductor is \(8\); for \(h=0\) there is no such \(\nu\).
Consequently the real primitive host associates are precisely
\(\rho\eta_q^i\eta_d^j\), with \(\rho=1\) or \(\nu\) when available
and \(i,j\in\{0,1\}\). Their conductors are
\(\operatorname{cond}(\rho)q^id^j\). The same product of conductors
holds if \(\rho\) is any primitive character induced by a character
of \(H\), since the three prime-power parts of \(Q\) are coprime.

For the two prime counts, use Thorner and Zaman's exceptional-zero
convention. For \(m\ge2\) the product of Dirichlet \(L\)-functions
modulo \(m\) has no zero in
\[
 \Re s\ge 1-\frac{1}{13\log(m(|\Im s|+3))}
\]
apart from a possible simple real zero \(\beta_m\) of a unique real
nonprincipal character \(\chi_m\bmod m\). If this zero exists, then
\(1-\beta_m\ge b_0 m^{-1/2}\) for an effective absolute \(b_0>0\).
\cite[arXiv v2, preceding Theorem~1.1 and (1.2)]{f1thorner-zaman}

For real \(U\ge8\) and \((a,m)=1\), let \(\lambda=1\) if there is
no exceptional zero. If there is one, let \(\lambda\) be the mean of
\(1-\chi_m(a)t^{\beta_m-1}\) on \((U/2,U]\), so that
\begin{equation}
\label{eq:family-prime-multiplier}
 \tfrac18\min\{1,(1-\beta_m)\log U\}<\lambda<2.
\end{equation}
Put \(\theta=32/37\) if the exceptional zero exists and \(7/12\)
otherwise. If \(\lambda U/(2\varphi(m))\ge U^{\theta+1/100}\), the
half-interval case of the relative prime number theorem is
\[
 \sum_{\substack{U/2<p\le U\\ p\equiv a\bmod m}}\log p
 =\frac{\lambda U}{2\varphi(m)}
 \left(1+O\left(\exp\left(
   -\frac{c\log U}{\log m+(\log U)^{2/5}(\log\log U)^{1/5}}
 \right)\right)\right).
\]
The constant \(c>0\) and the implied constant are effective and
uniform in \(U,m,a\).
\cite[arXiv v2, Theorem~1.1 and (1.7)]{f1thorner-zaman}

Use \((U/2,U]\subset(\sqrt U,U]\) for both candidate counts. For \(q\)
take \(m_0=8\), \(a=3\) and \(U=e^w\). These candidates satisfy \(q>w\)
for large \(Z\). With \(q\) fixed, the Chinese remainder theorem makes
the congruences for \(d\) a reduced class modulo
\[
  m_0=16q\prod_{\substack{3\le p\le w\\ p\ {\rm prime}}}p,
  \qquad U=\exp\bigl(w(\log Z)^2\bigr).
\]
For any real \(v>0\), Rosser and Schoenfeld give
\(\sum_{p\le v}\log p<1.01624v\)
\cite[Theorem~9, (3.32)]{rosser-schoenfeld}. At \(v=w\) this gives
\(\log m_0\le4\log U/(\log Z)^2\) for both candidate counts when \(Z\)
is large. Both choices also have \(\log U\ge Z^{\varepsilon/6}\).
The bounds on \(\beta_{m_0}\) and \(\lambda\) give
\(\lambda\gg m_0^{-1/2}\). Thus the main term is
\(\gg U m_0^{-3/2}=U^{1-o(1)}\), meeting the condition
\(U^{\theta+1/100}\) since \(\theta\le32/37\). The relative error
is \(o(1)\) uniformly in these moduli. Dividing by \(\log U\)
gives at least \(U^{1-o(1)}\) candidates, uniformly in the first
prime when choosing the second.

\Needspace{10\baselineskip}
\teamanchor{proof.auxiliary-primes}%
For each fixed \(J\ge1\) and all sufficiently large \(Z\), we need one
pair as above satisfying both assertions below. The classical proof for
\(M=1\) follows. For fixed \(M>1\), Alp\"oge's
\releasecite{Development/AutomorphicArithmetic/MeanSquareBounds/pairEx.lean}{112}{theorem on simultaneous quadratic twists}
supplies such a pair.

\begin{enumerate}
\item\label{item:centred-twists} For every character \(\xi\) modulo \(M\) and
  \((i,j)\in\{0,1\}^2\setminus\{(0,0)\}\), let \(\chi\) be the
  primitive character agreeing with \(\xi\eta_q^i\eta_d^j\) on the
  positive integers coprime to \(2Mqd\). For every real \(y\in[Z,Z^J]\),
  \[
    |G_\chi(y)|\le y^{1/2+\varepsilon/2}.
  \]
\item\label{item:twist-boxes} For every real primitive character \(\rho\) of conductor
  dividing \(8M\) and the same pairs \((i,j)\), the primitive
  character agreeing with \(\rho\eta_q^i\eta_d^j\) on those integers
  has no zero in
  \[
    \Re s>1-1/400,\qquad
    |\Im s|\le8M\exp\bigl(w+jw(\log Z)^2\bigr).
  \]
\end{enumerate}

Choose \(q\) first and then test twists by \(\eta_d\) in the progression
it determines. For any real \(R\ge3\),
Thorner and Zaman's density theorem gives
\[
 \sum_{f\le R}\ \sum_{\substack{\chi\bmod f\\ \chi\ \mathrm{primitive}}}
 \#\{\varrho:L(\varrho,\chi)=0,\ \Re\varrho>399/400,\
                              |\Im\varrho|\le R\}
 \le C R^{99/400}
\]
with an absolute effective \(C\). The count uses multiplicities and
includes real exceptional zeros.
\cite[arXiv v4, (1.1)--(1.2) and Theorem~1.2]{f1thorner-zaman-density}
In item~\ref{item:twist-boxes}, the two stages use \(R=8MU\) for \(q\)
and \(R=8Me^wU\) for \(d\). A failing twist has conductor divisible by
the candidate prime, with remaining factor at most \(8M\) or \(8Mq\),
respectively. That factor is smaller than every candidate at its
stage. Thus distinct rejected primes give distinct primitive
characters. Since \(R=U^{1+o(1)}\), these exclusions remove
\(o(U^{9/10})\) candidates.
For the general fixed target, the
\releasecite{Development/AutomorphicArithmetic/MeanSquareBounds/StagePow.lean}{41}{mean-square argument}
assumes every prime factor of the progression modulus is at most
\(w\) or greater than \(Z^{400J}\). The former holds for every factor of
\(8\) and for the small primes in the second modulus. Its other factor
\(q\) satisfies \(q>e^{w/2}>Z^{400J}\) for sufficiently large \(Z\).
The release controls the maximum on each dyadic block by a second
moment bound and uses it to choose one pair meeting both conditions.

For the characters in item~\ref{item:centred-twists}, evaluation at \(Z\) also gives
\[
 \left|\frac{L'(1,\chi)}{L(1,\chi)}\right|
 \le 1+ \frac{1}{Z}\sum_{n\le Z}\Lambda(n)\lfloor Z/n\rfloor
 =1+\frac{\log(\lfloor Z\rfloor!)}{Z}
 \le4\log y\qquad (y\ge Z).
\]

\paragraph{The zeta case from classical prime estimates.}
Fix \(J\ge1\), put \(\sigma=\tfrac12+\varepsilon/8\) and
\(T_0=Z^{2J}\).
\teamanchor{proof.real-character-zero-density}%
For \(Q_0,T_0\ge2\) and \(1/2\le\sigma<1\),
Heath-Brown's real-character density bound is
\[
 \sum_{\substack{\chi\ \mathrm{real,\ primitive}\\
                  \operatorname{cond}(\chi)\le Q_0}}
 N(\sigma,T_0,\chi)
 \ll_\alpha (Q_0T_0)^\alpha
      Q_0^{3(1-\sigma)/(2-\sigma)}T_0^{(3-2\sigma)/(2-\sigma)}
 \qquad(\alpha>0),
\]
where \(N(\sigma,T_0,\chi)\) counts zeros, with multiplicity, in
\(\sigma\le\Re s\le1\), \(|\Im s|\le T_0\)
\cite[Corollary~1 and Theorem~3]{f1heath-brown}.
For the first prime, take \(Q_0=U=e^w\) and test \(\eta_q\).
For the second, take \(U=\exp(w(\log Z)^2)\), \(Q_0=qU\) and test
\(\eta_d\) and \(\eta_q\eta_d\). The two single-prime characters are
odd and their product is even. Their conductors are \(q,d,qd\),
respectively, so each test assigns distinct real primitive
nonprincipal characters to distinct candidates at its stage.
In each stage \(Q_0=U^{1+o(1)}\) and \(T_0=U^{o(1)}\).
The conductor exponent is \(1-\delta\) for
\(\delta=(2\sigma-1)/(2-\sigma)>0\). Taking \(0<\alpha<\delta/4\)
excludes \(O(U^{1-\delta/2})\) candidates in each stage.
Intersect with the tall-rectangle exclusions above at \(M=1\).
They test \(\rho\eta_q\) first, and \(\rho\eta_d,\rho\eta_q\eta_d\)
second, for all real primitive \(\rho\) of conductor dividing \(8\).
Each failing character at a stage contains its candidate prime
in its conductor. The bounds \(O(U^{1-\delta/2})\) and
\(O(U^{99/400+o(1)})\) each exclude a proportion \(o(1)\) of the
\(U^{1-o(1)}\) candidates already counted in the progression.
We may choose \(q\) first and then \(d\) from the surviving primes,
with the second-stage bounds uniform in \(q\). This proves
item~\ref{item:twist-boxes} and leaves each of the three tested
characters without a zero in the Heath-Brown rectangle.

For a primitive nonprincipal character \(\chi\) of conductor \(f\),
let \(D_\chi=L'(1,\chi)/L(1,\chi)\) and
\(A_\chi(y)=\sum_{n\le y}\Lambda(n)\chi(n)(y/n-1)\).
The explicit formula over \(\mathbf Q\) gives, for real \(y\ge2\),
\[
 A_\chi(y)+(y-1)D_\chi
  =\sum_{\varrho}\frac{y^\varrho-1}{\varrho(1-\varrho)}
    +\frac{y-1}{2}\begin{cases}
         F_1(y), & \chi(-1)=1,\\
         F_3(y), & \chi(-1)=-1 .
       \end{cases}
\]
Here the zeros are the nontrivial zeros counted with multiplicity,
initially with a symmetric cutoff in height. The functions in the
formula are
\[
\begin{aligned}
 F_1(y)&=\log\frac{y+1}{y-1}
                +\frac{2}{y-1}\log\frac{y+1}{2},\\
 F_3(y)&=\log\frac{y^2}{y^2-1}
                +\frac{2}{y-1}\log\frac{2y}{y+1}.
\end{aligned}
\]
Indeed, at \(K=\mathbf Q\) the smoothed prime sum in the source
is \(A_\chi(y)/(y-1)\), and its constant for a nonprincipal
character is \(D_\chi\). The two archimedean choices follow from the
action on the real place. For \(y\ge2\) both archimedean terms are
\(O(\log(2y))\) with an absolute constant.
\cite[\S2.1 and Theorem~1, (2.2.1), pp.~256--257]{f1ihara-murty-shimura}
By the functional equation, the three selected real characters
have \(1-\sigma<\Re\varrho<\sigma\) for \(|\Im\varrho|\le T_0\)
\cite[(25.15.5)]{dlmf}.
Let \(N_\chi(t)\) count zeros with \(0<\Re\varrho<1\) and
\(|\Im\varrho|\le t\), including multiplicities. For a character
of conductor \(f>1\) and \(t\ge5/7\), Bennett--Martin--O'Bryant--Rechnitzer
give
\[
 \left|N_\chi(t)-\frac{t}{\pi}\log\frac{ft}{2\pi e}\right|
 \le \min\{0.247\log(ft)+6.894,\ 0.298\log(ft)+4.358\}.
\]
Apply this to the three primitive characters of conductors \(q,d,qd\).
\cite[Corollary~1.2, p.~1456]{f1bennett-zeros}
Thus the number of zeros in
\(n\le|\Im\varrho|<n+1\) is \(O(\log(f(n+3)))\). Use
\(N_\chi(n+1)-N_\chi(n-1)\) for \(n\ge2\), and \(N_\chi(2)\)
for \(n=0,1\).
For \(|\Im\varrho|\le T_0\) the reciprocal denominator is
\(O_\varepsilon((1+|\Im\varrho|^2)^{-1})\).
For \(|\Im\varrho|>T_0\) use \(|y^\varrho-1|\le2y\) and the
reciprocal bound \(|\Im\varrho|^{-2}\). Summing over the unit-height
strips gives an absolutely convergent zero sum and the bound
\[
 |A_\chi(y)+(y-1)D_\chi|
 \ll_\varepsilon y^\sigma\log(3f)
          +\frac{y}{T_0}\log(3fT_0)+\log(2y).
\]
The conductor and height satisfy
\(\log(3fT_0)\ll_J Z^{\varepsilon/5}(\log Z)^2\). For every
\(y\in[Z,Z^J]\) this gives
\[
 y^\sigma\log(3f)\ll_J y^{1/2+13\varepsilon/40}(\log y)^2,\qquad
  \frac{y}{T_0}\log(3fT_0)\ll_J Z^{-J+\varepsilon/5}(\log Z)^2.
\]
At \(y=Z\), use \(|A_\chi(Z)|\le\log(\lfloor Z\rfloor!)\) and
divide by \(Z-1\) to get \(|D_\chi|\le4\log Z\) for large \(Z\).
Adding \(D_\chi\) gives the centering in \eqref{eq:centred-prime-sums}.
Since \(13/40<1/2\), this proves
\(|G_\chi(y)|\le y^{1/2+\varepsilon/2}\) for all real
\(y\in[Z,Z^J]\) with the same pair, proving
item~\ref{item:centred-twists} for \(M=1\). For general \(M\), the twisted
characters in item~\ref{item:centred-twists} can be complex and
the release's prime-pair argument applies to those characters.

\paragraph{Mean squares in the two progressions.}
Let \(P=\{p:\sqrt U<p\le U,\ p\equiv a\bmod m_0\}\) be the
prime set at either stage above. For \(L=Z^A\), \(A\ge J\) fixed,
\(|\kappa(n)|\le1\) independent of \(p\), and \(y\in[Z,Z^J]\), set
\[
 S_L(y,p;\kappa)=\sum_{n\le L}\Lambda(n)\min(1,y/n)\kappa(n)
                         \Bigl(\frac np\Bigr).
\]
Then \(\sum_{p\in P}|S_L(y,p;\kappa)|^2\ll |P|y\log y\) with an
absolute constant for sufficiently large \(Z\).
For the relative factor \(|P|\), use the full-interval prime theorem
modulo \(m_0s\) at \(t\in[\sqrt U,U]\), for odd squarefree
\(1<s\le L^2\) coprime to \(m_0\).
Here \(\log(m_0L^2)\ll\log U/(\log Z)^2\), so
\(t\ge(m_0s)^{12}\) and the relative error is
\(O(\exp(-c(\log Z)^2))\) for some \(c>0\).
Summing over the reduced classes above \(a\) with weights
\((\frac{\cdot}{s})\) cancels the principal term. The secondary
term survives only if the exceptional character restricts to
\((\frac{\cdot}{s})\) on \((\mathbf Z/s\mathbf Z)^\times\).
The relative error is bounded using the sum of the nonnegative
main terms over those classes, including a possible bias from
a character induced from \(m_0\)
\cite[arXiv v2, Corollary~1.4 and Remark~1.5]{f1thorner-zaman}.
In the matching case the secondary term sums to zero without
the weights. Its signed contribution is negligible if
\(1-\beta\ge c/\log(m_0L^2)\), for fixed \(c>0\).
For \(Q_0>100\), the Landau--Page theorem gives an absolute
\(c_{\rm LP}>0\) such that among primitive characters of conductor
at most \(Q_0\), at most one character can have a real zero
\(\beta>1-c_{\rm LP}/\log Q_0\). An exceptional character is
quadratic \cite[Landau--Page; the real-axis case of Lemma~7.1 in the cited account]{f1page}.
Take \(Q_0=m_0L^2\) and \(c\le c_{\rm LP}\).
Since \(s\) is coprime to \(m_0\), the primitive quadratic
character of modulus \(s\) has conductor \(s\). In the matching
case the part of the exceptional conductor supported outside
\(m_0\) is therefore exactly \(s\). A single exceptional
character can affect at most one such \(s\).
Partial summation therefore gives
\(\sum_{p\in P}(\frac ps)=O_B(|P|Z^{-B})\) for every fixed \(B>0\),
with at most one exceptional \(s>1\).

Chebyshev's bound and partial summation make the contribution to
\(S_L\) from proper prime powers \(O(\sqrt y)\). Primes dividing
\(m_0\) contribute \(O(\log m_0)\).
For the remaining primes \(r\le L\), the weights
\(b_r=(\log r)\min(1,y/r)\) satisfy \(\sum_r b_r^2\ll y\log y\),
\(\max_r b_r\le\log y\) and \(\sum_r b_r\ll y(1+\log L)\).
Quadratic reciprocity on the fixed class modulo \(4\) reduces
the distinct-prime correlations to the last character average
with \(s\) the product of the two primes. At most one unordered
pair is exceptional. Its weight is \(O((\log y)^2)\), and the
other correlations contribute \(O_B(|P| Z^{-B} y^2(1+\log L)^2)\)
to the unnormalized moment. This gives the stated bound by
taking \(B>2J+1\).

For comparison, the unrestricted prime-modulus moment in
Montgomery--Vaughan is valid for arbitrary complex
\(a_1,\ldots,a_N\) and \(U>N^2\). It gives
\[
 \sum_{2<p\le U}\left|\sum_{n\le N}a_n\Bigl(\frac np\Bigr)\right|^2
 \ll\frac{U}{\log(2U/N^2)}
      \sum_{\substack{1\le s\le N\\s\ {\rm squarefree}}}
     \left(\sum_{m^2s\le N}|a_{m^2s}|\right)^2.
\]
With \(N=\lfloor L\rfloor\) our scales satisfy \(U>N^2\).
The preceding calculation provides the corresponding \(|P|\)
factor for the specified progression and the coefficients of \(S_L\).
\cite[Lemma~9, (7), p.~480]{f1montgomery-vaughan-moments}

For fixed \(M\) and any \(\tau>0\), the moment also gives
\(|G_\chi(y)|\le y^{1/2+\tau}\) on \([Z,Z^J]\) for a proportion
\(1-o(1)\) at each stage. Test the primitive associates of
\(\xi\eta_q\) at the first stage and of \(\xi\eta_q^i\eta_d\),
\(i=0,1\), at the second, for all \(\xi\bmod 8M\).
Apply the tall-rectangle exclusions to those associates. For a
surviving \(\chi\) of conductor \(f\), take \(L=Z^{3200J}\) and
\(\delta_0=1/400\).
We have \((\log f)^{10/\delta_0}\le L\le f^2\), and the zero-free
rectangle contains \(\Re s>1-\delta_0,\ |\Im s|\le L^{\delta_0}\).
The \(k=1\) case of Lamzouri--Languasco applies to any nonprincipal
character modulo a sufficiently large \(f\), when \(0<\delta_0<1/2\)
is fixed, \((\log f)^{10/\delta_0}\le L\le f^2\), and
\(L(s,\chi)\ne0\) throughout
\(\Re s>1-\delta_0,\ |\Im s|\le L^{\delta_0}\). It gives
\[
 -D_\chi=\sum_{n\le L}\frac{\Lambda(n)\chi(n)}{n}
                 +O_{\delta_0}(L^{-\delta_0/4}).
\]
Use \(f=\operatorname{cond}(\chi)\); the nonprincipal primitive
associate includes the candidate prime in its conductor, so
\(f\) is sufficiently large. The displayed bounds on \(L\) and
the tested rectangle supply the remaining hypotheses.
\cite[arXiv v1, Proposition~3.2, p.~8]{f1lamzouri-languasco}
Thus \(G_\chi(y)=-S_L(y,p;\kappa)+O(Z^{-J})\) with \(\kappa\)
the fixed factor before the new quadratic twist. Almost everywhere
in \(y\), \(y\,\partial_y S_L(y,p;\kappa)=
S_L(y,p;\kappa\mathbf1_{n>y})\).
On a block \([a_0,b_0]\subseteq[Z,Z^J]\), \(b_0\le2a_0\),
Cauchy--Schwarz bounds the squared maximum by twice the squared
value at \(a_0\) plus \(2a_0\) times the integral of the squared
derivative. The summed block maxima are \(O(|P|a_0\log(2a_0))\).
Markov's inequality and a dyadic sum give a discarded fraction
\(O(Z^{-2\tau_0}\log(2Z))\) for \(\tau_0=\min(\tau,1/4)\),
in addition to the rectangle exclusions. The finite union over
twists is uniform in \(q\) at the second stage. Evaluation at
\(Z\) with exponent \(1/2+\tau_0\) gives \(|D_\chi|\le4\log Z\).

\subsection{The Selmer bound on the same scales}

For a prime \(\ell\) let \(S_\ell\) be the places of \(E\) over
\(\ell\), and let \(G_{E,S_\ell}\) be the Galois group of the maximal
extension unramified outside \(S_\ell\). Write \(D_v\) for the
decomposition subgroup of a chosen prolongation of \(v\mid\ell\) to
an algebraic closure of \(E\). Restriction to \(D_v\) uses its map
to \(G_{E,S_\ell}\). For integers \(r\ge1\), put
\[
 s_\ell(E,r)=\operatorname{length}_{\mathbf Z_\ell}
  \ker\!\left(H^1(G_{E,S_\ell},(\mathbf Q_\ell/\mathbf Z_\ell)(-r))
     \longrightarrow\prod_{v\mid\ell}
         H^1(D_v,(\mathbf Q_\ell/\mathbf Z_\ell)(-r))\right).
\]
Here \(H^1\) is continuous Galois cohomology and the inverse \(r\)-th
power of the cyclotomic character acts on the coefficients. At finite
places \(v\), \(D_v\) identifies with the absolute Galois group of \(E_v\).

Let \(F\) be a totally imaginary number field and let \(S\) be a
finite set of finite places containing every place above \(\ell\).
For a finite \(\ell\)-primary \(G_{F,S}\)-module \(A\), write
\(\mathrm{Sha}^i_S(F,A)\) for the kernel of localization from
\(H^i(G_{F,S},A)\) to \(\prod_{v\in S}H^i(F_v,A)\).
Poitou--Tate duality gives a perfect pairing of finite groups
\[
 \mathrm{Sha}^1_S(F,A)\times
 \mathrm{Sha}^2_S\bigl(F,\operatorname{Hom}(A,\mu_{\ell^\infty})\bigr)
 \longrightarrow \mathbf Q/\mathbf Z .
\]
For this finite \(S\), the global cohomology groups and the local
targets with coefficients in \(A\) are finite.
The duality theorem includes the infinite places, whose local
cohomology in degrees \(1\) and \(2\) vanishes for such \(F\).
\cite[Theorem~I.4.10(a) and Corollary~I.4.15]{f1milne-adt}
Apply this with \(F=E\), \(S=S_\ell\) and
\(A=(\mathbf Z/\ell^k\mathbf Z)(-r)\), whose Tate dual is
\((\mathbf Z/\ell^k\mathbf Z)(r+1)\). Filtered direct limits commute
with the localization kernels, and finiteness at each level makes
the first derived inverse limits vanish. Passing to the direct and
inverse limits identifies the kernel defining \(s_\ell(E,r)\)
with the Pontryagin dual of the \(\ell\)-primary \'etale wild kernel
\[
 \ker\!\left(H^2(G_{E,S_\ell},\mathbf Z_\ell(r+1))
   \longrightarrow\prod_{v\mid\ell}H^2(E_v,\mathbf Z_\ell(r+1))\right).
\]

For a totally imaginary number field \(F\), every prime \(\ell\) and
integer \(r\ge1\), the \(K\)-theory comparison identifies
\[
 H^2_{\mathrm{\acute et}}(\mathcal O_F[1/\ell],\mathbf Z_\ell(r+1))
 \simeq K_{2r}(\mathcal O_F)[\ell^\infty],
\]
where \(K_{2r}(\mathcal O_F)\) is finite
\cite[Theorems~VI.8.2 and~VI.8.4]{f1weibel-kbook}. At finite coefficient
levels, the \'etale cohomology of \(\mathcal O_E[1/\ell]\) with locally
constant \(\ell\)-primary coefficients agrees with cohomology of its
\'etale fundamental group \(G_{E,S_\ell}\) because \(\ell\) is
invertible there \cite[Proposition~II.2.9]{f1milne-adt}. It follows in
the inverse limit that the wild kernel is contained in the
\(\ell\)-primary part of \(K_{2r}(\mathcal O_E)\). This includes the
case \(\ell=2\) since \(E\) is CM.
For fixed \(E,t\), the weighted sum over all primes is thus at most
\(\sum_{1\le r<t}\log|K_{2r}(\mathcal O_E)|\).
The required quantitative bound is
\begin{equation}
\label{eq:family-selmer}
  \sum_{\ell\in S}(\log\ell)\sum_{1\le r<t}s_\ell(E,r)\le B
  \qquad\text{for every finite set \(S\) of primes.}
\end{equation}
This inequality is uniform over the choices of algebraic closure
and prolongations.

\paragraph{The local representation data used in the count.}
\teamanchor{proof.selmer-galois-input}%
The release uses one semisimple Galois representation for a regular
algebraic polarized cuspidal pair \((\pi,\chi)\) on
\(\mathrm{GL}_a(\mathbf A_{F'})\), where \(F'\) is CM and may be an
auxiliary field in the counting argument. The polarization is
\(\pi^c\simeq\pi^\vee\otimes(\chi\circ N_{F'/F'^{+}}\circ\det)\), with
\(\chi_v(-1)=(-1)^a\) at the real places, allowing the replacement of
\(\chi\) by \(\chi\delta_{F'/F'^{+}}\) with unchanged pullback to \(F'\).
The normalized local component is
\(\pi_v\otimes|\det|_v^{(1-a)/2}\). With geometric reciprocity,
the inverse Artin map has valuation \(-\deg w\) on a Weil element \(w\)
of arithmetic degree \(\deg w\).
If \(\pi_\infty\) has the infinitesimal character of the dual of
highest weight \(\lambda_{\iota\tau}\), its coordinates at \(\tau\)
are \(z_j=-\lambda_{\iota\tau,a+1-j}+(a+1)/2-j\). Thus
\((a-1)/2-z_j\) gives the multiset
\(\{\lambda_{\iota\tau,j}+a-j\}_{j=1}^a\), using the convention
\(\mathrm{HT}(\mathbf Q_\ell(1))=\{-1\}\).
\releasecite{Development/AbelianVarietyTheory/PublishedTheorems/GaloisCohomology₂/GaloisRepOfPolarizedCuspidal.lean}{24}{Algebraic normalization}
\releasecite{Development/AbelianVarietyTheory/Proofs/InertiaEigenvalues/LocalArtinExistsHolds.lean}{19}{Local reciprocity convention}
\releasecite{Development/AbelianVarietyTheory/Definitions/GaloisGeometry/DualIsConjByOne.lean}{53}{Polarization}
\releasecite{Development/AutomorphicArithmetic/Foundations/GaloisRepresentations/TypeConditions/WeightsConjBis.lean}{54}{Weights and local requirements}.

For standard local Langlands, Barnet-Lamb, Gee, Geraghty and Taylor
give the semisimple representation with these labelled weights,
the conjugate multiset \(w-\mathrm{HT}_\tau\), and full
Frobenius-semisimplified compatibility away from \(\ell\).
Above \(\ell\) they give the de Rham property and
Frobenius-semisimplified compatibility in the Iwahori case, where
the representation is semistable, and crystalline in the spherical
case \cite[Theorem~2.1.1]{f1-blggt-pacw}. Caraiani's theorem gives the
full comparison above \(\ell\) for cohomological conjugate self-dual
cuspidal representations in rank at least two. The published statement
uses \(\operatorname{rec}(\Pi_v^\vee\otimes|\det|^{(1-a)/2})\) for
the parameter of its representation \(R_\ell(\Pi)\).
Take \(\Pi=\pi^\vee\). Its unramified Frobenius polynomials in the
construction of \(R_\ell(\pi^\vee)\) agree with those of
\(r_{\ell,\iota}(\pi)\), so semisimplicity and Chebotarev identify
the two representations \cite[Theorem~1.1]{f1-caraiani-lp}.
The condition
\(\pi^c\simeq\pi^\vee\) concerns the automorphic representation.
The specialization \(\psi=1\) concerns the analytic target and by itself
imposes no condition on the polarizing character of the representations
in the count.

For the essentially conjugate self-dual case, the result above \(\ell\)
up to semisimplification suffices for the standard parameter when that
parameter has zero monodromy
\cite[Theorem~A]{f1-blggt-lp2}. Indeed, apply the away-from-coefficient-prime
theorem at the same place with a different coefficient prime.
It shows that this standard parameter is pure. Since its monodromy
is zero, all its Frobenius eigenvalues have the same complex absolute
value. In the comparison above \(\ell\), the eigenvalues of the
semisimplified Weil representation are those same values.
The relation between the Weil action and monodromy sends each
generalized Frobenius eigenspace to one with absolute value rescaled
by the residue norm or its inverse. Hence the Galois parameter
also has zero monodromy, and the comparison is valid after
Frobenius semisimplification. This argument applies to the
standard parameter at ramified coefficient places in the
zero-monodromy range.

The release asks for comparisons with every correspondence satisfying
its stated axioms. For a toral parameter of zero monodromy, the
segment classification forces all segments to have length one.
The irreducible cuspidal parameters are then characters: an irreducible
complex Weil representation factoring through the abelianization is
one-dimensional. These parameters are determined by the common
Artin map. The second case of the release comparison allows a tame
parameter of zero monodromy at \(F'_v\simeq\mathbf Q_g\), with
\(g>16a^2\) or \(a=1\), and requires every repeated eigenvalue of a
tame inertia generator to have order dividing \(g-1\).
It asks for zero monodromy and the inertia characteristic polynomials.
For ramified \(v\mid\ell\) it also asks for potential semistability and
the labelled jumps in its local coefficient model. The comparison
for arbitrary correspondences in the second case uses the
depth-zero type law, including the inertia polynomial of an inflated
cuspidal type. We use the release's construction for this comparison
and for its transport to the specified local models. The existence of
a correspondence with that type law is a separate input of the release.
\releasecite{Development/AutomorphicArithmetic/Foundations/GaloisRepresentations/TypeConditions/WeightsConjBis.lean}{26}{Tame and toral comparisons}
\releasecite{Development/AutomorphicArithmetic/Foundations/LocalLanglands/IsRecTOfIsRec.lean}{28}{Depth-zero type law}
\releasecite{Development/AutomorphicArithmetic/Foundations/LocalLanglands/LocalCentralUnits/RecTExists.lean}{20}{Existence with the specified type law}
\releasecite{Development/AutomorphicArithmetic/Foundations/GaloisRepresentations/IntegralStructures/RowSrc.lean}{6}{Galois representation on the stated models}.

\paragraph{Special values for the critical parity.}
\teamanchor{proof.shared-critical-selmer-values}%
Let \(L/\mathbf Q\) be abelian CM of degree \(n_L\), with maximal
real subfield \(F_L\), and let \(N_L\) be divisible by the conductors
of its primitive Dirichlet characters. For \(k\ge2\) put
\[
 H_{\ell,k}(L)=H^2_{\mathrm{\acute et}}
          (\mathcal O_L[1/\ell],\mathbf Z_\ell(k)),
 \qquad w_k(L)=|H^0(L,\mathbf Q/\mathbf Z(k))|.
\]
Use the corresponding notation for \(F_L\). The factor
\(\Lambda_k(L)\) is the product of \(L(1-k,\xi)\) over primitive
associates of the characters of \(L\) with
\(\xi(-1)=(-1)^k\). By Artin factorization it equals
\(\zeta_{F_L}(1-k)\) for even \(k\), and
\(L(1-k,\chi_{L/F_L})\) for odd \(k\).

For even \(k\), Huber--Kings at odd primes and Kolster at \(2\)
give
\[
 \log_\ell|H_{\ell,k}(F_L)|
    =v_\ell(\Lambda_k(L))+v_\ell(w_k(F_L)).
\]
The real regulator in this parity has rank zero. The torsion
in degree one is the \(\ell\)-part of \(w_k(F_L)\) by the
coefficient sequence. The Euler factors removed above \(\ell\)
are \(\ell\)-units at \(1-k\).
\cite[Theorems~1.3.1 and~1.4.1]{f2-huber-kings}
\cite[Appendix~A, Theorem~A.1 and Corollary~A.3]{f2-rognes-weibel-dyadic}

For odd \(k\ge3\), Kolster's higher relative class number formula
for abelian CM fields gives
\[
 \log_\ell\frac{|H_{\ell,k}(L)|}{|H_{\ell,k}(F_L)|}
 =v_\ell(\Lambda_k(L))+v_\ell(w_k(L))
    +v_\ell(Q_k)-([F_L:\mathbf Q]+1)v_\ell(2),
 \qquad Q_k\in\{1,2\}.
\]
These are the ordinary cohomology groups of the respective rings
of \(\ell\)-integers. In the formulation of Kolster--Sands
apply Theorem~6.1 to the cyclic extension \(L/F_L\) of degree
two, with \(L\) abelian over \(\mathbf Q\) and \(k=2j+1\).
Its \(h^-_k(L)\) is the quotient of the products over all primes
of the orders of these groups for \(L\) and \(F_L\), respectively.
The theorem evaluates \(L(1-k,\chi_{L/F_L})\) as
\((-1)^{j[F_L:\mathbf Q]}2^{[F_L:\mathbf Q]+1}
h^-_k(L)/(Q_k w_k(L))\). Taking \(v_\ell\) gives the displayed
identity with the indicated dyadic term. Here \(Q_k\) is the
index in \(H^1_{\mathrm{\acute et}}(L,\mathbf Z_2(k))\) of the
sum of the restricted group from \(F_L\) and the torsion subgroup
\(H^0(L,\mathbf Q_2/\mathbf Z_2(k))\).
\cite{f2-kolster-relative}\cite[\S\S2 and~5, Theorem~6.1]{f2-kolster-sands}
\cite[Propositions~1.2--1.3 and p.~225]{f2-voight-relative}

Let \(P_k=\prod_{p\le k+1}p\), and let \(f_\xi\) be the primitive
conductor of \(\xi\). The special-value formula is
\[
 L(1-k,\xi)=-\frac{f_\xi^{\,k-1}}{k}
      \sum_{a=1}^{f_\xi}\xi(a)B_k(a/f_\xi).
\]
It follows by decomposing the Dirichlet series into Hurwitz zeta
functions \cite[25.11.14 and~25.11.36]{dlmf}.
Expanding the Bernoulli polynomial and using von Staudt--Clausen
on \(B_j\) for \(2\le j\le k\) shows that
\(kP_k f_\xi L(1-k,\xi)\) is an algebraic integer.
The same expansion and \(|B_j|\le j!\) give
\(|B_k(u)|\le3k!\) on \([0,1]\), so
\(|L(1-k,\xi)|\le3(k f_\xi)^k\).
The parity list is Galois-stable and the factors are nonzero
by the functional equation. Thus
\[
 M_k(L)=\Lambda_k(L)\prod_{\xi(-1)=(-1)^k} kP_k f_\xi
\]
is a nonzero rational integer, and \(P_k\le(2k)^k\) gives
\[
 \log|M_k(L)|
 \le n_L\bigl(\log3+\log(kN_L)+k\log(2k)+k\log(kN_L)\bigr).
\]
In particular \(v_\ell(\Lambda_k(L))\le v_\ell(M_k(L))\).

If \(\ell^a\mid w_k(L)\), the cyclotomic image in
\((\mathbf Z/\ell^a\mathbf Z)^\times\) has index at most \(n_L\)
and exponent dividing \(k\). Its order is at most \(k\) at odd
\(\ell\) and at most \(2k\) at \(2\), by the cyclic structure of
the odd prime-power units and the two cyclic factors at \(2\).
It follows that \(\ell^a\le2n_Lk\) at odd \(\ell\) and
\(2^a\le4n_Lk\) at \(2\). Each odd prime in \(w_k(L)\) contributes
at most \(\log(2n_Lk)\) to its logarithm, and there are at most
\(n_Lk\) of them. In particular
\(\log w_k(L)\ll n_Lk\log(2n_Lk)\).

\paragraph{Critical parity at odd primes on the classical host.}
\teamanchor{proof.selmer-minus-parts}%
For the remainder of the arithmetic comparison take \(\psi=1\), write
\(F=E^{\langle c_\infty\rangle}\), \(n_E=[E:\mathbf Q]\), and
\(Q=2^{h+2}qd\). Put \(k=r+1\) for a twist \(1\le r<t\).
For odd \(\ell\), the release calls the \((-1)^r\)- and
\((-1)^{r+1}\)-eigenspaces of the fine Selmer group the plus and minus
parts, respectively. The Poitou--Tate pairing is equivariant for
complex conjugation, so the latter corresponds to the
\((-1)^k\)-eigenspace of the wild kernel. Let \(s_\ell^-(E,r)\) be its
length and write \(H^i(L)=H^i(\mathcal O_L[1/\ell],\mathbf Z_\ell(k))\).
By \eqref{eq:classical-host-inertia}, \(E/F\) is unramified at every
finite place, and the finite \'etale double cover gives
\(H^i(F)\simeq H^i(E)^{c_\infty=1}\) by transfer at odd \(\ell\).
\releasecite{Development/AutomorphicArithmetic/AutomorphicVarieties/LengthEqAddEigen.lean}{68}{The sign of each Selmer component}.

Use the
\hyperlink{proof.shared-critical-selmer-values}{critical special values}
with \(L=E\) and \(N_L=Q\). For even \(k\), the required invariant
part of \(H^2(E)\) has order \(|H^2(F)|\). For odd \(k\),
the anti-invariant part has order \(|H^2(E)|/|H^2(F)|\).
These are the two odd-prime transfer consequences above.
Since \(w_k(F)\mid w_k(E)\), both formulas yield
\[
 s_\ell^-(E,r)\le v_\ell(\Lambda_k(E))+v_\ell(w_k(E))
               \le v_\ell(M_k(E))+v_\ell(w_k(E)).
\]
Summing over odd primes with weights \(\log\ell\), the bounds for
\(M_k(E)\) and the odd part of \(w_k(E)\) give, for \(t\ge2\),
\begin{equation}
\label{eq:family-critical-selmer}
 \sum_{\ell\ {\rm odd}}\log\ell \sum_{1\le r<t}s_\ell^-(E,r)
       \le 5n_Et^2\log(2n_EtQ).
\end{equation}

\paragraph{Fixed exceptional towers.}
\teamanchor{proof.selmer-cyclotomic-ranks}%
For a fixed abelian field containing \(\mu_\ell\), or \(\mu_4\) when
\(\ell=2\), adjoining the \(\ell\)-power roots gives its cyclotomic
\(\mathbf Z_\ell\)-extension. The Ferrero--Washington theorem
gives \(\mu_\ell(X)=0\) for the norm limit \(X\) of the
\(\ell\)-primary ideal class groups in the cyclotomic
\(\mathbf Z_\ell\)-extension of a field abelian over \(\mathbf Q\).
The structure theorem for finitely generated torsion
\(\mathbf Z_\ell[[\Gamma]]\)-modules then makes \(X\) finitely
generated over \(\mathbf Z_\ell\). Its quotient obtained by
allowing \(S_\ell\)-integers is finitely generated as well.
\cite[main theorem, p.~377]{f2-ferrero-washington}
The class group at each sufficiently high layer
is a quotient of this module on its \(\ell\)-primary part. Indeed,
above the first layer where a prime is totally ramified in the
remaining tower, the Hilbert \(\ell\)-class field of the layer
intersects the tower trivially over that layer. Thus the
\(\ell\)-ranks are bounded for the fixed base field
\cite{f2-ferrero-washington}. A bound uniform in moving base fields
requires a separate argument.

For the classical target the fixed exceptional set is \(\{3,7\}\)
and the fixed fields are \(K_\ell=\mathbf Q(\zeta_\ell)\) for
\(\ell=3,7\). In its definition \(p_K=7\), \(p_1=3\), and the
class number of \(\mathbf Q(\sqrt{-7})\) is one by Minkowski's bound.
\releasecite{Development/AutomorphicArithmetic/AutomorphicVarieties/StmtFields.lean}{57}{Exceptional primes}
\releasecite{Development/AutomorphicArithmetic/AutomorphicVarieties/PrimeSetBoundDataOfPrimitiveChar.lean}{71}{Fixed auxiliary fields}.
The Minkowski bounds for \(K_3\) and \(K_7\) are less than \(2\)
and \(5\), respectively. In \(K_7\), \(2\) and \(3\) have residue
degrees \(3\) and \(6\), so every class has an ideal representative
of norm one. Consequently both class numbers are one.
In each cyclotomic tower the unique prime above \(\ell\) is totally
ramified at every step, with all other primes unramified.
Iwasawa's one-prime criterion states that in a finite cyclic
\(\ell\)-extension with just one ramified prime and that prime
totally ramified, the vanishing of the \(\ell\)-primary class
group at the base implies its vanishing at the top.
Here the base \(K_\ell\) is totally imaginary and its class
number is one, so applying the criterion to each layer over
\(K_\ell\) gives trivial \(\ell\)-primary class groups throughout
the tower \cite[pp.~257--258]{f2-iwasawa-class}. The class groups of
the rings of \(\ell\)-integers are quotients of those groups, so
their inverse-limit \(\ell\)-class module is zero. The
wild-kernel codescent isomorphism is as follows. For a number
field \(K\), assume \(\ell\) is odd or \(\mu_4\subset K\)
when \(\ell=2\). Let \(K'=K(\mu_{2\ell})\), let \(K'_\infty\)
be its cyclotomic \(\mathbf Z_\ell\)-extension, and put
\(\Gamma=\operatorname{Gal}(K'_\infty/K')\),
\(\Delta=\operatorname{Gal}(K'/K)\). For every \(r\ge1\) the
localization kernel in \(H^2(G_{K,S_\ell},\mathbf Z_\ell(r+1))\)
is canonically \((X'(r))_{\Gamma\times\Delta}\), where \(X'\)
is the norm limit of the \(\ell\)-primary class groups of the
rings of \(\ell\)-integers in \(K'_\infty\).
\cite[Theorem~3.1, p.~446]{f2-validire}
Apply this at each layer as \(K\). These fields already contain
\(\mu_\ell\), the remaining tower has \(X'=0\), and \(\Delta=1\).
This supplies the
fixed-field layer bound with \(B=0\) in the character count.
\releasecite{Development/AutomorphicArithmetic/AutomorphicVarieties/HeightAndVanishing/LayerLe.lean}{12}{Fixed-layer input and aggregate bounds}.
The moving host's \(3\)- and \(7\)-primary plus parts still enter
the estimate on \(E\).

\paragraph{A bound for the full even \(K\)-group.}
\teamanchor{proof.selmer-lemma-k}%
Let \(F\) be totally imaginary of degree \(d_F\) and absolute
discriminant \(D_F\). Write \(\Delta_F=D_F^{1/d_F}\) and put
\(j=2e\ge2\). We will show, with an absolute implied constant, that
\begin{equation}
\label{eq:family-even-k-group}
 \log\bigl(2+\log|K_j(\mathcal O_F)|\bigr)
    \ll (jd_F)^5\log(2jd_F\Delta_F).
\end{equation}

By the plus construction and homological stability over Dedekind
domains there is a Hurewicz homomorphism
\[
 K_j(\mathcal O_F)\longrightarrow H_j(\Gamma,\mathbf Z),
 \qquad \Gamma=\mathrm{GL}_N(\mathcal O_F),\quad N=4j+5 .
\]
In this range inclusion into the infinite general linear group
induces an isomorphism in degree \(j\)
\cite[Theorem~3.1(i)]{f1charney-dedekind-stability}.
Arlettaz's integral Hurewicz theorem gives an annihilator \(R_j\)
for the kernel of this very map. In his notation
\(R_j=\prod_{k=2}^j L_k\), where each \(L_k\) divides the product of
the primes \(p\le k/2+1\)
\cite[Definitions~7.11 and~7.14, Corollary~7.19(a)]{f1arlettaz-k-theory}.
In particular the kernel is killed by
\[
 a_j=\left(\prod_{p\le j/2+1,\ p\text{ prime}}p\right)^{j-1},
 \qquad \log a_j\ll j^2\log(2j).
\]

Let \(\Gamma(3)\) be the principal subgroup modulo \(3\mathcal O_F\),
of index \(I\le3^{d_FN^2}\). It is torsion-free. In an integral basis
of \(\mathcal O_F^N\) its matrices lie in the principal subgroup
modulo \(3\) of \(\mathrm{GL}_{d_FN}(\mathbf Z)\).
A nontrivial finite-order element would have a power of prime
order \(p\) of the form \(1+3^aB\), where \(a\ge1\) and \(B\) is
nonzero modulo \(3\). Its \(p\)-th power is nonidentity modulo
\(3^{a+1}\) for \(p\ne3\) and modulo \(3^{a+2}\) for \(p=3\), by
the binomial formula. Compose the Hurewicz map with the integral
homology transfer to \(\Gamma(3)\). Transfer followed by inclusion
is multiplication by \(I\), so the kernel \(B_j\) of the composite
is killed by \(b=Ia_j\).

We bound the image using Ash's contractible well-rounded retract.
The congruence subgroup acts freely since all cell stabilizers are
finite. The first barycentric subdivision gives a quotient complex
computing \(H_*(\Gamma(3),\mathbf Z)\)
\cite[\S2.3]{f1bayer-fluckiger-emery-houriet}
\cite[\S4.3]{f1ash-mcconnell-retract}.
In the retract, the minimal vectors \(M(\sigma)\) on the interior
of a cell \(\sigma\) determine it. Passing from a cell to a face
enlarges its set of minimal vectors. Write \(n=d_FN\). In each
well-rounded lattice the minimal vectors span over \(F\), so there
is an \(\mathcal O_F\)-basis whose vectors have \(h\)-length at most \(B\), with
\(\log(2B)\ll Nd_F^2\log(2Nd_F\Delta_F)\)
\cite[Proposition~4.7]{f1bayer-fluckiger-emery-houriet}.
In these basis coordinates take one complex embedding from each
conjugate pair and use ordinary Euclidean volume on \(\mathbf R^n\).
The real matrix \(Q\) of \(q_h\) has
\(\operatorname{tr} Q=2\sum_i q_h(e_i)\le2NB^2\), and the covolume
of \(\mathcal O_F^N\) is at most \(D_F^{N/2}\).
The open ellipsoid \(q_h<1\) contains no nonzero lattice vector.
Minkowski's convex body theorem and the cube of side \(2/\sqrt n\)
inside the Euclidean unit ball give
\(\det Q\ge n^{-n}D_F^{-N}\). The least eigenvalue is at least
\(\det Q/(\operatorname{tr} Q)^{n-1}\), so every minimal vector has
Euclidean norm at most \(R\), where
\[
 R^2=n^n D_F^N (2NB^2)^{n-1}.
\]
The basis bound permits the same \(B\) and \(R\) for every cell.
The points of the standard lattice \(\mathcal O_F^N\) have pairwise
distance at least one, since a nonzero algebraic integer has
absolute norm at least one. Take \(\Phi\) to be its points in the
closed Euclidean ball of radius \(R\). Packing disjoint open balls
of radius \(1/2\) shows that
\[
 \log(1+|\Phi|)\le \log\bigl(1+(1+2R)^n\bigr)
 \ll N^3d_F^4\log(2Nd_F\Delta_F).
\]
Thus one \(\Gamma\)-translate of each minimal set lies in the same
\(\Phi\).
Let \(m\le d_FN(N+1)/2\) be the dimension of the real space of
Hermitian forms. For any flag of cells, place the minimal vectors of
its smallest cell in \(\Phi\). All the other minimal sets in the
flag are subsets of this set. In order of increasing minimal set,
extend a basis of the linear functionals \(h\mapsto q_h(v)\) by
choosing vectors from each set when its span increases. A basis for
a set of minimal vectors determines its cell. Indeed all the
corresponding equations \(q_h(v)=1\) follow from the basis
equations because the coefficients of any such linear relation
sum to one, by evaluation on a form in the cell. If two cells
have the same basis, these equations give both inclusions of
their exact minimal sets. Thus every strict step in a flag
increases the rank. The ordered list of chosen vectors, with
the end of each basis marked, determines the entire flag.
The number of \(\Gamma\)-orbits of subdivided simplices is at most
\((1+2|\Phi|)^m\), by padding each list to length \(m\).

The quotient by \(\Gamma(3)\) has at most
\(A=I(1+2|\Phi|)^m\) simplices. Each boundary column from degree
\(j+1\) has \(\ell^1\)-norm at most \(j+2\), counting faces before
any identifications. The torsion in degree \(j\) is the torsion of
the cokernel of this boundary since the image of the next boundary
is free abelian. The determinant bound for the torsion of a
cokernel gives \(\log|H_j(\Gamma(3),\mathbf Z)_{\rm tors}|\le
A\log(j+2)\)
\cite[Lemma~3.1]{f1bayer-fluckiger-emery-houriet}.
The bound on \(|\Phi|\), together with \(N=O(j)\), gives
\[
 \log\bigl(2+A\log(j+2)\bigr)
 \ll N^5d_F^5\log(2Nd_F)+N^5d_F^4\log D_F
 \ll (jd_F)^5\log(2jd_F\Delta_F).
\]

For the kernel, write \(A_p=K_{2e}(\mathcal O_F)[p^\infty]\).
The even \(K\)-group is finite and \(A_p\) is
\(H^2(\mathcal O_F[1/p],\mathbf Z_p(e+1))\)
\cite[VI, Theorems~8.2 and~8.4]{f1weibel-kbook}. Put \(F'=F(\mu_p)\), of degree \(d'\le d_Fp\).
Its degree over \(F\) divides \(p-1\), so coefficient reduction
and restriction inject \(A_p/pA_p\) into
\(H^2(\mathcal O_{F'}[1/p],\mu_p)\). The reduced twists become
isomorphic to \(\mu_p\) over \(F'\). The extension \(F'/F\)
is unramified away from \(p\) and tame above \(p\), giving
\(D_{F'}^{1/d'}\le\Delta_F p\) by the relative different formula.
Kummer theory bounds the order of the target by \(h_{F'}p^{d'}\).
The Brauer invariants are supported on the places over \(p\), since
\(F'\) is totally imaginary
\cite[VI, (8.1.1)]{f1weibel-kbook}.
Minkowski gives an ideal of norm at most \(2^{d'}\sqrt{D_{F'}}\)
in every class, and the number of ideals of norm at most \(Y\)
is at most \(Y^2\zeta_{F'}(2)\le Y^22^{d'}\). Hence
\[
 \log|A_p/pA_p|\ll d_Fp\log(2p\Delta_F).
\]
The same computation at \(p=2\) has \(F'=F\). By the elementary
divisors and \(bB_j=0\) we now have
\[
 \log|B_j|
 \le\sum_{p\mid b}v_p(b)\log|A_p/pA_p|
 \ll d_F b\log(2b\Delta_F)\log b .
\]
Here \(\log b\ll d_Fj^2+j^2\log(2j)\). Taking
\(\log(2+\log|B_j|)\) gives a bound of the size in
\eqref{eq:family-even-k-group}. The kernel and the image together
give \eqref{eq:family-even-k-group}.

For each prime \(\ell\), put \(R=\mathcal O_F[1/\ell]\). In
particular \(H^2(R,\mathbf Z_\ell(e+1))\) is finite, and
\begin{equation}
\label{eq:family-lemma-k}
 \log|H^2(R,\mathbf Z_\ell(e+1))|
 \le \exp\bigl(C(2ed_F)^5\log(4ed_F\Delta_F)\bigr)
\end{equation}
with an absolute \(C\). Finite-coefficient Galois and \'etale
cohomology agree for this ring of \(\ell\)-integers and these twists.
Finiteness of \(H^1\) at each finite level gives the comparison
in the inverse limit.
\releasecite{Formalism/GaloisCohomology/Statement.lean}{54}{Localizations and Tate twists}.
For \(u=2ed_F\ge4\), an absolute \(C_K\) satisfies
\(C u^5\log(2u\Delta_F)\le u^{C_K}\log(2D_F)\).
This is the cohomological estimate called Lemma~K. The release uses it in a
comparison on a bounded-degree dyadic layer and in its odd-prime
count. The classical all-layer dyadic argument below uses parity
formulas.
\releasecite{Development/AutomorphicArithmetic/AutomorphicVarieties/GeneratorCountTuples/LemmaK.lean}{9}{The cohomological estimate}
\releasecite{Development/AutomorphicArithmetic/AutomorphicVarieties/LayeredIsomorphisms/PartA.lean}{24}{The dyadic comparison}
\releasecite{Development/AutomorphicArithmetic/AutomorphicVarieties/HeightAndVanishing/CountVs.lean}{86}{The exceptional odd-prime application}
\releasecite{Development/AutomorphicArithmetic/AutomorphicVarieties/ModularityArgument/ThresholdArguments/DenseThresholdCases/WindowBoundOfLiteratureLevel.lean}{25}{Use in the automorphic counts}.

\paragraph{The dyadic bound at every layer of the classical family.}
\teamanchor{proof.selmer-classical-two-window}%
\teamanchor{proof.selmer-classical-all-layers}%
Let \(q\equiv3\bmod8\) and \(d\equiv7\bmod16\) be distinct primes.
Put \(B_j=\mathbf Q(\zeta_{2^{j+2}})^+\),
\(E_j=B_j(\sqrt{-q},\sqrt{-d})\), \(L_j=E_j(i)\), and
\(F_j=L_j^+=B_j(\sqrt q,\sqrt d)\). For every \(j\ge0\) and \(r\ge1\),
write \(N=2^{j+2}qd\), \(n=[L_j:\mathbf Q]=8\cdot2^j\), and \(k=r+1\).
With an absolute implied constant, we have
\begin{equation}
\label{eq:classical-dyadic-all-layers}
 s_2(E_j,r)\ll n k\log(2nkN).
\end{equation}

Indeed, let \(c\) be complex conjugation. Mizusawa treats
\(\mathbf Q(\sqrt{-qd})(\sqrt q,\sqrt d)=L_0\) and its cyclotomic
\(\mathbf Z_2\)-extension. His Lemma~4.2 gives the inverse limit of
the ordinary \(2\)-class groups of the \(L_j\) as a free
\(\mathbf Z_2\)-module of rank three. Its proof gives odd ordinary
class number for every \(F_j\)
\cite[\S4.2, Lemma~4.2 and its proof, p.~126]{f2-mizusawa-dyadic}.
The operator \(1+c\) on each \(2\)-class group factors through the
class norm to \(F_j\), so it vanishes. These conclusions give at
most three generators and \(c=-1\) on the inverse limit \(X'\) for
the rings of \(2\)-integers. Indeed it is a quotient of the ordinary
limit by exactness of inverse limits of finite groups.

Write \(L=L_j,\ F=F_j\). For \(K=F,L\) let
\(H_K=H^2_{\mathrm{\acute et}}(\mathcal O_K[1/2],\mathbf Z_2(k))\),
which is finite for \(k\ge2\)
\cite[VI, \S8]{f1weibel-kbook}\cite{f2-kolster-relative}.
Let \(A\) be the kernel of localization of \(H_L\) at the dyadic
places. The codescent isomorphism identifies \(A\) with the
coinvariants of \(X'(r)\) under \(\operatorname{Gal}(L_\infty/L)\),
since \(i\in L\) \cite[Theorem~3.1]{f2-validire}. On the dual this
is the group of invariant continuous homomorphisms from \(X'\) to
\((\mathbf Q_2/\mathbf Z_2)(-r)\), with \(c\) acting on both source
and target. Indeed, for the global cyclotomic group
\(\operatorname{Gal}(L_\infty/L)\) and each dyadic decomposition
subgroup, a generator minus one acts on these divisible coefficients
by a nonzero scalar, giving \(H^1=0\), and \(H^2=0\) by cohomological
dimension one. Inflation--restriction identifies the strict
kernel with the invariant homomorphisms unramified everywhere and
trivial at the dyadic places. Thus \(A\) has at most three generators and \(c\)
acts on it by \((-1)^k\).
By Poitou--Tate duality \(s_2(L,r)=\log_2|A|\)
\cite[Theorem~I.4.10(a)]{f1milne-adt}.
The extension \(L=F(i)\) is quadratic and unramified away from \(2\)
at finite places, so restriction and corestriction are available on the
two rings of \(2\)-integers.

If \(k\) is even, then \(2A=(1+c)A\) lies in the restriction of
\(H_F\), and hence \(\log_2|A|\le3+\log_2|H_F|\).
For \(L=L_j\) use the
\hyperlink{proof.shared-critical-selmer-values}{critical special values}
with \(N_L=N\). For even \(k\) this evaluates the term
\(\log_2|H_F|\) as \(v_2(\Lambda_k(L))+v_2(w_k(F))\).

Suppose \(k\) is odd, so \(k\ge3\). Restriction \(H_F\to H_L\)
has kernel killed by two and image in \(H_L^c\).
The finite group \(H_F\) satisfies
\(\log_2|H_F[2]|=\log_2|H_F/2|\le2n\). Indeed reduction modulo two
embeds \(H_F/2\) in \(H^2(\mathcal O_F[1/2],\mu_2)\).
The Picard group has odd order, and the \(2\)-torsion of the Brauer
group is supported at the dyadic and real places by the Kummer
sequence and the local invariant sequence.
The invariant subgroup \(A^c=A[2]\) has order at most eight.
The quotient \(H_L^c/A^c\) embeds in \(H_L/A\), which embeds
in the product of the local \(H^2\)'s at \(2\).
For a dyadic completion of degree \(b\), local duality bounds the
length of this group by \(\log_2(4br)\). In fact its dual consists
of the invariants in \((\mathbf Q_2/\mathbf Z_2)(-r)\).
If their order is at least \(2^a\), the cyclotomic image modulo
\(2^a\) has index at most \(b\) and exponent dividing \(r\).
The \(r\)-torsion in \((\mathbf Z/2^a\mathbf Z)^\times\) has at
most \(2r\) elements, so \(2^a\le4br\). Thus
\(\log_2|H_F|\le 2n+3+n\log_2(4nr)\).

For odd \(k\) the same formula for the relative class number,
together with the displayed bound for \(\log_2|H_F|\), bounds
\(\log_2|H_L|\) and hence \(\log_2|A|\). In both parities use
\(v_2(\Lambda_k(L))\le v_2(M_k(L))\) and
\(v_2(w_k(F))\le v_2(w_k(L))\le\log_2(4nk)\). Thus
\(\log_2|A|\ll nk\log(2nkN)\).
Restriction of the strict fine Selmer group from \(E_j\) to \(L\)
has kernel of order at most two by inflation--restriction, since
the coefficient invariants over \(L\) form a finite cyclic group.
This proves \eqref{eq:classical-dyadic-all-layers}.

Fix \(0<\varepsilon<1/52\), an admissible profile \(h\) and
\(A_0\ge1\). Whenever the pair is chosen at \(Z=X^{1/J}\), for
every \(x\in[Z,X]\) we have \(n\le16x^{2+8\varepsilon}\) at
\(j=h(x)\), and
\(\log N\ll\log(2x)+x^{\varepsilon/5}(\log x)^2\).
Summing \eqref{eq:classical-dyadic-all-layers} for
\(1\le r<\lfloor x^{A_0}\rfloor\) thus gives a polynomial bound
in \(x\), at most \(\exp(x^{4\varepsilon/5})\) for all sufficiently
large \(x\), uniformly in \(q,d\). The dyadic assertion alone
allows \(J=2\) and \(W=\sqrt X\). On any larger fixed ambient
range the estimate is available on every subinterval needed for
the other inputs.

\paragraph{The odd-primary contribution from smaller layers and twists.}
\teamanchor{proof.selmer-classical-low-layer}%
At odd \(\ell\), the character idempotents of
\(\operatorname{Gal}(E_h/\mathbf Q)\) split the strict group over
an unramified coefficient ring \(O\) containing all the character
values. The sum of the \(O\)-lengths is the original
\(\mathbf Z_\ell\)-length, by extension of scalars. The summands in
the plus part are those with \(\rho(c_\infty)=(-1)^r\). A character
\(\rho\) has a least layer \(s=s(\rho)\) through which it factors.
Its summand has the same length on \(E_h\) and \(E_s\). In fact
restriction identifies the strict group on \(E_s\) with the
\(\operatorname{Gal}(E_h/E_s)\)-invariants of the strict group on
\(E_h\). Inflation--restriction gives the isomorphism for absolute
cohomology and injectivity on the local groups for inertia away
from \(\ell\) and decomposition at \(\ell\). In each case the quotient
group has order dividing \(2^{h-s}\), a unit at \(\ell\).

Write \(n_s=[E_s:\mathbf Q]=4\cdot2^s\) and \(Q_s=2^{s+2}qd\).
The conductor--discriminant formula gives
\(|\operatorname{disc} E_s|^{1/n_s}\le Q_s\). By Poitou--Tate
duality, the strict group at each \(\ell\) is dual to a subgroup
of the \(\ell\)-primary part of \(K_{2r}(\mathcal O_{E_s})\).
Its length multiplied by \(\log\ell\) is the logarithm of its
order. Summing over the distinct odd primary parts and using
\eqref{eq:family-even-k-group} gives
\[
 \log\!\left(2+\sum_{\ell\ {\rm odd}}(\log\ell)s_\ell(E_s,r)\right)
 \ll (r n_s)^5\log(2r n_s Q_s).
\]

Fix \(J\) and choose a pair at \(Z=X^{1/J}\). At every
\(x\in[Z,X]\), put \(h=h(x)\). Counting each \(\rho\) at its least
layer, the plus summands with \(r n_{s(\rho)}\le x^{\varepsilon/8}\)
have total weighted length \(o(\exp(x^\varepsilon/2))\), uniformly
over the pairs. Indeed \(h=O(\log x)\) and
\(\log Q_s\le\log Q_h=O(x^{\varepsilon/5}(\log x)^2)\) for \(s\le h\).
The logarithm of the bound for each eligible \((s,r)\) is
\(O(x^{33\varepsilon/40}(\log x)^2)\), and there are at most
\((h+1)x^{\varepsilon/8}\) such pairs. Descent to \(E_s\) allows
each group of summands to be bounded by the whole group on that
layer. This argument includes the moving fields and the
coefficient primes \(3,7,q,d\).

\paragraph{The remaining classical Selmer bound.}
\teamanchor{proof.selmer-classical-residual}%
Two further sets of plus summands vanish. At \(3,7\), large \(q,d\)
give \(\ell\nmid Q_h\). The part of \(\rho\) prime to \(\ell\) away
from the \(\ell\)-conductor is therefore \(\rho\) itself. The
odd-order condition of Case~IV forces \(\rho=1\) on this \(2\)-group
host. For that character the length is the strict length over
\(\mathbf Q\) by descent, and restriction to
\(\mathbf Q(\zeta_\ell)\) injects into the zero group proved above.
For \(\rho=1\) at an odd prime with \(\ell-1\mid 2r\), the plus
condition forces \(r\) even. The residual coefficient character
\(\bar\kappa_\ell^{-r}\), for cyclotomic character \(\kappa_\ell\),
is trivial or cuts out the real quadratic field
\(K=\mathbf Q(\sqrt{\ell})\).
In the trivial case, a class killed by \(\ell\) lifts to an
additive character with values in \(\mathbf F_\ell\), by the
multiplication-by-\(\ell\) coefficient sequence. It vanishes on
inertia away from \(\ell\) and therefore factors through
\(\operatorname{Gal}(\mathbf Q(\mu_{\ell^\infty})/\mathbf Q)\)
by Kronecker--Weber \cite[Ch.~I, Theorem~4.16]{f1milne-cft}.
The decomposition group at \(\ell\) surjects onto this quotient,
through which the divisible coefficient
action also factors. A local coboundary at \(\ell\) is therefore
a global coboundary. In the quadratic case, the residual
character is nontrivial on \(\ell\)-inertia, so the divisible
coefficients have no invariants there. A class killed by
\(\ell\) lifts to a residual cocycle with zero localization at
\(\ell\) and zero inertia away from \(\ell\). Restriction to
\(K\) is injective and gives an unramified character of its class
group. In this field every class has an ideal representative of
norm at most \(\sqrt\ell\). The number of such ideals is at most
\(\sum_{m\le\sqrt\ell}d(m)\le\sqrt\ell(1+\log\sqrt\ell)<\ell\).
In both cases the \(\ell\)-torsion of the strict group vanishes,
which gives full vanishing. These arguments account for the
classical specializations of Counts IV and Z.

\teamanchor{proof.selmer-dyadic-quadratic}%
The pure dyadic quadratic character \(\nu=(2/\cdot)\) also has zero
strict \(3\)-primary summand for every \(r\ge1\). Its fixed-factor
field is \(F_8=\mathbf Q(\sqrt2)\). Set
\(L_j=F_8(\zeta_{3^{j+1}})\) for \(j\ge0\).
The conductor--discriminant formula for the three quadratic
characters gives \(|\operatorname{disc}(L_0)|=8\cdot3\cdot24=576\).
The Minkowski bound is \(36/\pi^2<4\). There are no ideals of norm
\(2\), since \(2\) is inert in \(\mathbf Q(\zeta_3)\), or of norm
\(3\), since \(3\) is inert in \(F_8\). Hence \(L_0\) has class
number one. The unique prime over \(3\) in \(F_8\) has unramified
quadratic completion over \(\mathbf Q_3\). Adjoining the cyclotomic
fields, whose completions at \(3\) are totally ramified, gives a
unique prime above \(3\) on each \(L_j\), totally ramified over
\(L_0\). All other primes are unramified in \(L_j/L_0\). The
one-prime criterion gives \(\mathrm{Cl}(L_j)[3^\infty]=0\) for every
\(j\) \cite{f2-iwasawa-class}. Therefore the inverse limit of the
\(3\)-parts of the class groups of \(\mathcal O_{L_j}[1/3]\) is
zero. Schneider's codescent gives zero wild kernel with coefficients
\(\mathbf Z_3(r+1)\) on \(L_0\), which already contains \(\mu_3\).
The Poitou--Tate pairing above gives zero strict group with
coefficients \((\mathbf Q_3/\mathbf Z_3)(-r)\).
\cite[Theorem~3.1]{f2-validire}

For a finite Galois extension of degree prime to \(3\), restriction
identifies the strict groups downstairs with the invariant strict
groups upstairs. Use inflation--restriction in absolute cohomology
and the description of the local conditions by inertia away from
\(3\) and full decomposition groups at \(3\). The local group
quotients have order dividing the extension degree, including at
ramified primes, so local restriction is injective. The degree-two
extension \(L_0/F_8\) thus gives vanishing on \(F_8\). For \(h\ge1\),
the extension \(E_h/F_8\) has degree \(2^{h+1}\), and the
\(\nu\)-summand is contained in the invariant group for this extension.
This proves the claimed vanishing on \(E_h\). Since \(\nu\) is even,
its plus summand occurs for even \(r\). These terms at \(3\) belong
to Count~II, since \(3\) is exceptional and the character away from
the \(3\)-conductor has order two.

\paragraph{Pure dyadic characters at odd primes.}
\teamanchor{proof.selmer-pure-dyadic}%
Let \(B_s=\mathbf Q(\zeta_{2^{s+2}})^+\), with \(B_0=\mathbf Q\).
Write \(\mathcal S_\ell(L,r)\) for the strict group with coefficients
\(A_r=(\mathbf Q_\ell/\mathbf Z_\ell)(-r)\) over a number field \(L\).
A pure dyadic character \(\theta\) of the classical host factors
through \(\operatorname{Gal}(B_h/\mathbf Q)\). If its least layer is
\(s\), restriction identifies its strict summand on \(E_h\) with the
one on \(B_s\). In fact \(E_h/B_s\) has degree \(4\cdot2^{h-s}\).
Inflation--restriction in absolute cohomology identifies the
downstairs group with the upstairs invariants. It also gives
injective restriction on inertia cohomology away from \(\ell\) and
on decomposition cohomology at \(\ell\), since the relevant local
quotient orders divide the extension degree. The coefficients are
unramified away from \(\ell\), so these are exactly the strict
conditions. This argument also applies to any finite Galois extension
of degree prime to \(\ell\). For the character projections, one can
extend scalars to a finite unramified coefficient ring containing the
host character values. Those values have order a power of \(2\),
so the idempotents are integral and the lengths are unchanged
after this scalar extension.

Suppose \(\ell\equiv3,5\bmod8\) and \(\ell-1\mid r\).
The representative \(\ell^*\in\{\ell,-\ell\}\) with
\(\ell^*\equiv1\bmod4\) is \(5\bmod8\).
Successive squaring gives
\(v_2((\ell^*)^{2^j}-1)=2+j\) for \(j\ge0\).
The Frobenius of \(\ell\) therefore generates
\(\operatorname{Gal}(B_s/\mathbf Q)\), and \(B_s\) has a single place
\(v\) above \(\ell\), with unramified completion over \(\mathbf Q_\ell\).
The local cyclotomic extension is totally ramified, so the local and
global cyclotomic images on \(B_s\) are both \(\mathbf Z_\ell^\times\).
In particular \(A_r^{G_{B_s}}=A_r^{G_{(B_s)_v}}\).

The coefficient sequence for multiplication by \(\ell\) lifts any
strict class killed by \(\ell\) to
\(H^1(G_{B_s,S_\ell},A_r[\ell])\).
The localization of the lift lies in the image of the local connecting
map from \(A_r^{G_{(B_s)_v}}/\ell A_r^{G_{(B_s)_v}}\). Subtract the
global connecting class of the same invariant. The new lift has zero
localization at \(v\). The residual coefficient \(A_r[\ell]\) is
trivial because \(\ell-1\mid r\), so the lift is an everywhere
unramified additive character. At a real place an additive character
of odd order is trivial. Global class field theory thus puts the lift
in \(\operatorname{Hom}(\operatorname{Cl}(B_s),A_r[\ell])\).
Horie's Proposition~3 for the cyclotomic \(\mathbf Z_2\)-extension
of \(\mathbf Q\) asserts vanishing of its \(\ell\)-class group
for every prime \(\ell\equiv3,5\bmod8\). In the notation of the
proposition and Lemma~1, the class group is the direct limit
under extension of ideals, and the maps on \(\ell\)-parts are
injective because the relative degrees are powers of \(2\).
Our \(B_s=\mathbf Q(\zeta_{2^{s+2}})^+\) is its degree-\(2^s\)
layer, so \(\operatorname{Cl}(B_s)[\ell^\infty]=0\).
\cite[Lemma~1, p.~260, and Proposition~3, p.~289]{f2-horie-dyadic-class}
Hence the strict group has zero
\(\ell\)-torsion. Every element is of \(\ell\)-power order, so the whole
group vanishes. All pure dyadic plus summands at \(3\) follow by taking
every positive even \(r\). The same proof applies at every
\(\ell\equiv3,5\bmod8\) to the indicated multiples of \(\ell-1\).

\paragraph{Vanishing above a fixed dyadic layer.}
\teamanchor{proof.selmer-pure-dyadic-high}%
Let \(s\ge1\) and let \(\theta\) have exact order \(2^s\).
For even \(r\) put \(K_\ell=\mathbf Q(\zeta_\ell)^+\) and
\(L_s=B_sK_\ell\). One can take \(K_\ell=\mathbf Q\) instead when
\(\ell-1\mid r\). The fields \(B_s\) and \(K_\ell\) are disjoint by
their relatively prime cyclotomic conductors. The extension \(L_s/B_s\)
has degree dividing \(\ell-1\), so strict restriction is injective
and is equivariant for the \(B_s\)-factor of
\(\operatorname{Gal}(L_s/\mathbf Q)\). On \(G_{L_s}\) the
residual coefficient \(A_r[\ell]\) is trivial, since the cyclotomic
image modulo \(\ell\) over \(K_\ell=\mathbf Q(\zeta_\ell)^+\) is
\(\{1,-1\}\); the alternative \(K_\ell=\mathbf Q\) uses
\(\ell-1\mid r\).

Write \(H_s=\operatorname{Gal}(L_s/K_\ell)\) and let
\(D\subset H_s\) stabilize a place \(w\mid\ell\).
Under the identification with \(\operatorname{Gal}(B_s/\mathbf Q)\)
it is generated by \(\operatorname{Frob}_\ell\), since \(K_\ell\) is
totally ramified at \(\ell\) and \(B_s\) is unramified there.
Every element of \(D\) admits a lift on the local algebraic closure
fixing \(\mathbf Q_\ell(\mu_{\ell^\infty})\). Indeed it acts on the
unramified completion of \(B_s\), which is disjoint from the totally
ramified cyclotomic extension. Thus \(D\) acts trivially on
\(A_r^{G_{(L_s)_w}}\) and on the image of the local connecting map.
The product of these images over \(w\mid\ell\) is the permutation
representation on \(H_s/D\) with that trivial coefficient. If
\(\theta(\operatorname{Frob}_\ell)\ne1\), its \(\theta\)-part is zero.
This holds also for the residual idempotents because reduction is
injective on roots of unity of order prime to \(\ell\).

Lift an element of the strict \(\theta\)-summand killed by \(\ell\) to
residual cohomology on \(L_s\) and project the lift with the
\(\theta\)-idempotent. Every localization above \(\ell\) is zero by
the preceding calculation. The lift is therefore a class group
character in the \(\theta\)-part. The odd-primary part of
\(\operatorname{Cl}(L_{s-1})\) at \(\ell\) injects into that of
\(\operatorname{Cl}(L_s)\) by extension of ideals, since its norm
composite is multiplication by \(2\). If their orders are equal,
extension is surjective and the relative involution acts trivially
on the class group at \(\ell\). The involution acts by \(-1\) in
the \(\theta\)-part. The lift is then zero, proving vanishing of the
entire strict \(\theta\)-summand.

Washington's theorem states that if \(F/\mathbf Q\) is abelian,
\(F_\infty/F\) is the cyclotomic \(\mathbf Z_p\)-extension, and
\(\ell\ne p\) is prime, the sequence
\(v_\ell(|\operatorname{Cl}(F_s)|)\) for its degree-\(p^s\)
layers is eventually constant
\cite[main theorem, p.~87]{f2-washington-non-p}.
Take \(p=2\) and \(F=K_\ell\). The two fields have coprime
cyclotomic conductors at each finite dyadic layer, so \(L_s\)
is the degree-\(2^s\) layer over \(K_\ell\). Thus the required
class group orders agree for \(s\) sufficiently large at each
fixed odd \(\ell\). The other condition holds when
\(s>v_2(\ell^*-1)-2\), since the image of \(\ell^*\) in the
cyclic group of order \(2^s\) is then nontrivial and \(\theta\) is
faithful. This yields a threshold \(s_\ell\) such that every pure
dyadic \(\theta\) of least layer \(s>s_\ell\) has zero strict
\(\ell\)-summand at every even twist. For \(\ell-1\mid r\), the
theorem of Fukuda and Komatsu gives a cutoff also for a varying
\(\ell\). Write \(h_s=|\operatorname{Cl}(B_s)|\). Put
\(c_\ell=v_2(\ell-1)\) for \(\ell\equiv1\bmod4\) and
\(c_\ell=v_2(\ell^2-1)\) for \(\ell\equiv3\bmod4\).
Their theorem gives \(\ell\nmid h_s/h_{m_\ell}\) for every
\(s\ge m_\ell=2c_\ell-3+\lfloor\log_2\ell\rfloor\).
Thus \(s>m_\ell\) puts both successive layers in the stationary
range when \(K_\ell=\mathbf Q\). Their Corollary~1.1 and
Theorem~1.2 give \(\ell\nmid h_s\) for every \(s\ge1\) if
\(\ell<1.2\cdot10^8\) or \(\ell\equiv\pm9\bmod16\).
\cite[Theorems~1.1--1.2 and Corollary~1.1, pp.~360--361]{f2-fukuda-komatsu-weber}
These criteria are used with \(s>v_2(\ell^*-1)-2\).

\paragraph{A common interval for the pure dyadic terms at \(3,7\).}
\teamanchor{proof.selmer-pure-dyadic-fixed-primes}%
Take \(\mathcal P_{\rm fix}=\{3,7\}\) before choosing the window
parameters. By the preceding vanishing and prime-to-\(\ell\)
restriction, for an even \(r\) the total length of pure dyadic
summands on any \(E_h\) at \(\ell\in\mathcal P_{\rm fix}\) is at most
\(\operatorname{length}\mathcal S_\ell(M_\ell,r)\), where
\(M_\ell=B_{s_\ell}(\zeta_\ell)\). Indeed the pure summands on
\(E_h\) are the group on \(B_h\); for \(h>s_\ell\) its surviving
characters all factor through \(B_{s_\ell}\). For \(h\le s_\ell\)
restriction into \(B_{s_\ell}\) is injective, and the remaining
restriction to \(M_\ell\) has degree prime to \(\ell\).

Set \(M_{\ell,\infty}=B_{s_\ell}(\mu_{\ell^\infty})\),
let \(\Gamma_\ell=\operatorname{Gal}(M_{\ell,\infty}/M_\ell)\), and
choose its generator \(\gamma_\ell\) with cyclotomic image
\(a_\ell=1+\ell\). The conductor of \(B_{s_\ell}\) is a power of \(2\),
so this image generates the full \(1+\ell\mathbf Z_\ell\).
Let \(X'_\ell\) be the inverse limit, under norms, of the
\(\ell\)-primary class groups of rings of \(\ell\)-integers in this
tower. By Ferrero--Washington it is a finitely generated
\(\mathbf Z_\ell\)-module. Poitou--Tate and Schneider's codescent
identify the strict group's dual with
\((X'_\ell(r))_{\Gamma_\ell}=X'_\ell/(\gamma_\ell-a_\ell^{-r})X'_\ell\).
Here \(M_\ell\) is abelian and already contains \(\mu_\ell\), so there
is no auxiliary prime-to-\(\ell\) coinvariant in codescent.
\cite{f2-ferrero-washington}
\cite[Theorem~3.1]{f2-validire}
Let \(D\) be the sum over \(\ell\in\mathcal P_{\rm fix}\) of the ranks of the
free quotients of \(X'_\ell\). Their characteristic polynomials
for \(\gamma_\ell\) have a total of \(D\) roots \(\beta\), with
multiplicities. The torsion submodules have fixed finite lengths.
The strict length is bounded by the torsion length plus the valuation
of the determinant on the free quotient. This determinant is nonzero
for \(r\ge1\) because the strict group is finite. Its valuation is
the sum of \(v_\ell(a_\ell^{-r}-\beta)\) over the corresponding roots.

For any \(X>1\), each root gives at most one integer
\(1\le r\le X^{400}\) with
\(v_\ell(a_\ell^{-r}-\beta)>2+400\log_\ell X\). For distinct
\(r,r'\) in that range,
\[
 v_\ell(a_\ell^{-r}-a_\ell^{-r'})=1+v_\ell(r-r')
 \le 1+400\log_\ell X.
\]
This proves the claim by the ultrametric inequality. Call the resulting
at most \(D\) integer weights exceptional for this \(X\).
Set \(\delta=\varepsilon/30\). On the fixed fields \(M_\ell\),
\eqref{eq:family-lemma-k} and the wild-kernel inclusion bound the
total weighted length from \(r\le x^\delta\) by
\[
 O_{\mathcal P_{\rm fix}}\!\left(
 x^\delta\exp\bigl(C_{\mathcal P_{\rm fix}}
 x^{5\delta}\log(2x)\bigr)\right).
\]
Here \(5\delta<\varepsilon\).
For \(x^\delta<r<\lfloor x^{400}\rfloor\), when no exceptional
weight is in this range, the contribution with weights \(\log\ell\)
is \(O_{\mathcal P_{\rm fix}}(x^{400}\log X)\). Both estimates are
\(o(\exp(x^\varepsilon/2))\) uniformly for \(Z\le x\le X\) with
\(Z=X^{1/J}\) and \(J\) fixed.

Choose an integer \(J\ge2\) before \(X,q,d\) with
\[
 \log J>D\log(400/\delta)+(D+1)\log(3/2).
\]
For each exceptional \(r\ge2\), remove the closed real interval
\([r^{1/400},r^{1/\delta}]\) from \([Z,X]\). In the variable
\(\log\log x\) each removed interval has length at most
\(\log(400/\delta)\). At most \(D+1\) complementary intervals
remain, and one has length greater than \(\log(3/2)\). It contains
a closed subinterval corresponding to \([U,U^{3/2}]\).
Thus on this interval the pure dyadic mass at all primes in
\(\mathcal P_{\rm fix}\) is \(o(\exp(x^\varepsilon/2))\) at the unchanged
cutoff \(\lfloor x^{400}\rfloor\). The choice is independent of
\(q,d\) and the estimates hold for every admissible pair and
profile there.

Put \(T_x=\lfloor x^{400}\rfloor\) and \(B(x)=\exp(x^\varepsilon/2)\).
Let \(A(x)\) be the total mass with weights \(\log\ell\) through
\(1\le r<T_x\) from the \(2\)-primary terms, the odd minus terms,
the small-layer and small-twist range proved above, the vanishing
summands, and the pure dyadic summands at \(3,7\), counting each
term once. The bound \(A(x)=o(B(x))\) holds uniformly on the
selected interval \([U,U^{3/2}]\). The bounds
\eqref{eq:family-critical-selmer} and
\eqref{eq:classical-dyadic-all-layers} contribute polynomial terms
there since \([E_h:\mathbf Q]=O(x^{2+8\varepsilon})\).
The pure dyadic terms at \(3,7\) satisfy the bound on that very
interval by the preceding calculation.

For the remaining odd plus summands it suffices that their weighted
sum at each point of this same interval be at most \(B(x)-A(x)\),
uniformly over finite prime sets. A bound \(B(x)/2\) suffices at large
scales. Adding the two parts proves \eqref{eq:family-selmer} at \(T_x\)
throughout \([U,U^{3/2}]\), with finiteness from the wild-kernel inclusion.
Conjugating the local decomposition groups leaves the kernel unchanged.

The release supplies a pointwise bound on the plus terms throughout
\([Z,X]\). With the fixed auxiliary layer terms at \(3,7\) set to zero,
it gives \((Q_h T_x)^C\) with \(C\) fixed before the host. This applies
to the remaining terms on the same \(q,d,h(x)\) throughout the
chosen subinterval. The count uses \(2T_x+1<q,d\) and fixed lower
thresholds on both primes. These hold eventually since \(J\) is
fixed and \(\log q>\tfrac12 Z^{\varepsilon/5}\). Since
\(\log(Q_h T_x)=O(x^{\varepsilon/5}(\log x)^2)\) for the actual modulus,
the bound is \(o(B(x))\) throughout \([Z,X]\) and hence on the
selected interval.
\releasecite{Development/AutomorphicArithmetic/AutomorphicVarieties/HeightAndVanishing/DirichletPrimeToLPartMulLPowerPart.lean}{164}{Host modulus and size conditions}
\releasecite{Development/AutomorphicArithmetic/AutomorphicVarieties/HeightAndVanishing/LayerLe.lean}{39}{The pointwise plus estimate}
\releasecite{Development/AutomorphicArithmetic/AutomorphicVarieties/HeightAndVanishing/JointWinOfCounts.lean}{16}{Assembly of the counts}.
\rhrange{On these fields the release bounds the weighted plus mass by
\((Q_hT_x)^C=o(B(x))\) throughout \([Z,X]\).}
{First choose the interval for the pure dyadic terms at \(3,7\).
Use the source estimate for the remaining odd plus terms on each
field of that same interval.}
Case~III, for nontrivial odd-order host characters, is empty.
The bounded-order branch of Case~V, called Count~Vs in the release,
lies in the proved small-layer range at twist one for large scales.
Independent control of the remaining terms requires the aggregate
in Case~I, the character estimates in Case~II, and the remaining
Case~V estimates at \(3,7\) and twist one. Case~I includes the
coefficient primes \(q,d\) since \(2r<q-1,d-1\). At \(3,7\),
the terms of Case~II outside the proved ranges involve a \(q\)- or
\(d\)-factor. Case~V at twist one uses odd characters.
\releasecite{Development/AutomorphicArithmetic/AutomorphicVarieties/HeightAndVanishing/Cases.lean}{8}{Partition of the character cases}
\releasecite{Development/AutomorphicArithmetic/AutomorphicVarieties/HeightAndVanishing/CountVs.lean}{86}{The bounded-order part at twist one}.

\teamanchor{proof.admissible-host-windows}%
Fix a nonnegative integer \(h(y)\) for each real \(y\ge1\) with
\(y^{2+8\varepsilon}\le 2^{h(y)}\le2y^{2+8\varepsilon}\).
For \(\psi=1\) we may take
\(h(y)=\lceil(2+8\varepsilon)\log_2 y\rceil\).
Fix the integer \(J\) from the determinant comparison before selecting
\(X\) and the prime pair. Take \(X\) large enough that \(Z=X^{1/J}\)
exceeds the fixed thresholds. The classical prime theorem selects
\(q,d\) at \(Z\) with the analytic bounds on \([Z,X]\). The determinant
comparison gives the interval \([U,U^{3/2}]\) on which to combine
the pure dyadic estimate with the remaining pointwise Selmer bound.
Every real \(x\) of this interval uses the same \(q,d\) and has
\eqref{eq:family-selmer} at \(\lfloor x^{400}\rfloor\) on
\(E(h(x),q,d)\). The inequalities
\(x^{1/J}\le Z\le x\le X=Z^J\) identify the prime-selection data
used in the host. We can meet any lower endpoint request by
taking \(X\) large enough. Fix a sequence of these intervals with
each left endpoint beyond the preceding right endpoint. Their union
is a single set of real scales with arbitrarily distant full windows
and a specified \(Z,q,d\) on each interval. The field and bound are
established before the choice of ranks, levels and translates.

For the Dirichlet extension, use the
\hyperlink{candidate.family}{general Selmer window assertion}
with its larger twist cutoff at the rescaled argument below. The
release derives that theorem from local--global compatibility and
base change, and supplies the compatibility from theta lifts.
\releasecite{Development/AutomorphicArithmetic/AutomorphicVarieties/ModularityArgument/ThresholdArguments/DenseThresholdCases/VwinOfRow.lean}{4}{The general Selmer window}
\releasecite{Development/AutomorphicArithmetic/CharacterBoxes/MainConclusion.lean}{7}{Its application}
\releasecite{Development/AutomorphicArithmetic/Foundations/GaloisRepresentations/IntegralStructures/RowSrc.lean}{6}{The required local--global compatibility}.

For \(M>1\), use the \(J\) of the window assertion when
\releasecite{Development/AutomorphicArithmetic/CharacterBoxes/WinPairs.lean}{34}{choosing the pair \(q,d\) at \(Z=X^{1/J}\)}.
Take \(X\) large enough for the prime estimates at that scale. The
release's prime-pair theorem supplies it for a general fixed target.
The resulting Selmer interval lies in \([Z,X]=[Z,Z^J]\), on which
the prime-pair bounds hold. The subgroup
\[
 H_0=\operatorname{Gal}\!\left(
   \mathbf Q(\zeta_{2^{h+2}})^+\mathbf Q(\psi)/\mathbf Q(\psi)\right)
 \ \le H
\]
is cyclic, and its index \(o=[H:H_0]\) is the order of \(\psi\).
The intersection of the real dyadic field with \(\mathbf Q(\psi)\)
stabilizes for large \(h\), since \(\mathbf Q(\psi)\) has finite degree.
Consequently \(n=|H|=c_\psi2^h\) for a constant \(c_\psi\ge1\).
Put \(k_0=c_\psi^{1/(2+8\varepsilon)}\) and \(Y=k_0W\). If
\(W\ge k_0^3\), then \(Y^{3/2}\le W^2\), and both \(x\) and
\(y=x/k_0\) lie in \([W,W^2]\) for every \(x\in[Y,Y^{3/2}]\).
In particular \(x^{1/J}\le Z\le x\), as required for the host
construction. For sufficiently large \(X\) all the layers \(h(y)\)
lie in the stabilized range. Choose the dyadic layer \(h(y)\) at \(x\).
Then
\[
 x^{2+8\varepsilon}\le n\le2x^{2+8\varepsilon}.
\]
Taking \(W\ge k_0^{400}\) gives \(y^{401}\ge x^{400}\), and
\(y\le x\).
The Selmer bound therefore holds at cutoff \(\lfloor x^{400}\rfloor\)
with \(B=\sqrt b\), where \(b=\exp(x^\varepsilon)\).
Taking \(X\) arbitrarily large gives arbitrarily distant intervals
with \(q,d\) fixed throughout each interval. The dyadic layer can vary
with \(x\). For \(\psi=1\) we have \(o=c_\psi=1\) and \(H_0=H\).

\subsection{Level primes, sampling sets and character estimates}

\paragraph{Sizes and cutoff for the classical host.}
\teamanchor{proof.host-realization}%
At every sufficiently large \(x\) in the classical intervals, the
field above satisfies
\begin{equation}
\label{eq:classical-host-scale}
 n=|H|=2^{h(x)},\qquad Q_S=4n,\qquad
 x^{2+8\varepsilon}\le n\le2x^{2+8\varepsilon},\qquad
 f=[F:\mathbf Q]=2n .
\end{equation}
The prime-pair bounds and \(x^{1/J}\le Z\le x\) give
\[
 \tfrac12 x^{\varepsilon/(5J)}<\log q\le x^{\varepsilon/5},
 \qquad
 \log d\le x^{\varepsilon/5}(\log x)^2 .
\]
Also \(\log d>\log q\) once \((\log Z)^2>2\). Hence \(d>q>x^A\)
eventually for every fixed \(A\).
As \(\log Q_S=O(\log x)\), these upper bounds give
\(\log(Q_Sqd)\le2Z^{\varepsilon/5}(\log Z)^2\) for large \(Z\):
\(\log Q_S=O_J(\log Z)\), \(\log q\le Z^{\varepsilon/5}\), and
\(\log d\le Z^{\varepsilon/5}(\log Z)^2\). As \(Z\le x\), this
also gives \(\log(Q_Sqd)\le2x^{\varepsilon/5}(\log x)^2\).
In particular every fixed power of \(Q_Sqd\) is
\(\exp(o(x^\varepsilon))\).
For the root discriminant, put \(K=\mathbf Q(\zeta_{Q_Sqd})\) and
write \(X^{Q_Sqd}-1=\Phi_{Q_Sqd}(X)R(X)\) with \(R\) integral.
At \(\zeta_{Q_Sqd}\) the product of the derivative of \(\Phi_{Q_Sqd}\)
and \(R\) equals \(Q_Sqd\) times a root of unity. The nonzero norm of
\(R(\zeta_{Q_Sqd})\) is an integer, and the norm of the derivative
is the discriminant of the power basis in absolute value.
The field discriminant divides it, so \(\operatorname{rd}(K)\le Q_Sqd\).
The discriminant formula for the tower \(K/E/\mathbf Q\) gives
\(\operatorname{rd}(E)\le\operatorname{rd}(K)\). Thus
\(\log\operatorname{rd}(E)=o(x^\varepsilon)\).

The prescribed rank relations are
\[
 x^{1/2-\varepsilon/16}\le m\le2x^{1/2-\varepsilon/16},\qquad
 x^{9/2+10\varepsilon}\le k\le2x^{9/2+10\varepsilon},\qquad
 r=k+m-1,\quad m+3\le k .
\]
The integer \(t\) used as twist cutoff is specified by
\begin{equation}
\label{eq:classical-rank-cutoff}
\begin{aligned}
 r+t &=8fr+4+\lceil x^{377}\rceil\,A_x,\\
 A_x &=8\left\lceil\frac{10r}{8}
     \left(f+\frac{f^2(k-1)^3x^{2\varepsilon}}{m}\right)\right\rceil .
\end{aligned}
\end{equation}

\paragraph{Choice of ranks on the classical field.}
\teamanchor{proof.classical-rank-choice}%
On the frame in \eqref{eq:classical-host-scale} let
\[
 m=2\left\lceil x^{1/2-\varepsilon/16}/2\right\rceil,\qquad
 k=2\left\lceil x^{9/2+10\varepsilon}/2\right\rceil,\qquad r=k+m-1 .
\]
For \(A\ge2\) the even integer \(2\lceil A/2\rceil\) lies in
\([A,2A]\). At all sufficiently large \(x\) we have \(m\ge2\),
\(m+3\le k\) and \(m\le 2n\le k\), since the exponents for
\(m,n,k\) are strictly increasing in that order. Take \(t\) from
\eqref{eq:classical-rank-cutoff} with \(f=2n\). The integer \(r\)
is odd and the right side of that formula is divisible by \(4\),
so \(t\) is odd. Also \(t\ge7r+4\), whence \(r+5\le t\),
\(3m<t\) and \(m^2\le k+m(t+1)\). This gives all the rank
conditions with \(n\) and the Artin map already specified by the
classical field. The cutoff estimate following
\eqref{eq:classical-rank-cutoff} therefore applies on this field.
\releasecite{Formalism/LocalRingCardinality/Params.lean}{6}{The rank conditions}
\releasecite{Development/AbelianVarietyTheory/Proofs/GlobalInterpolation/ExistsRanksRankValFrame.lean}{45}{The released construction on a prescribed frame}.
For \(x\ge1\) and \(0<\varepsilon<1/52\) we have
\(f\le4x^{2.16}\), \(k\le2x^{4.70}\), \(r\le4x^{4.70}\),
\(m\ge x^{0.49}\), and \(x^{2\varepsilon}\le x^{0.04}\).
Consequently \(f^2(k-1)^3x^{2\varepsilon}/m\le128x^{17.97}\),
so \(A_x\le5288x^{22.67}\). Using
\(\lceil x^{377}\rceil\le2x^{377}\) in
\eqref{eq:classical-rank-cutoff} gives \(t\le10708x^{399.67}\).
For \(x\ge2^{100}\) the last bound is at most \(x^{400}\).
Since \(t\) is an integer with \(t\le\lfloor x^{400}\rfloor\),
the classical arithmetic window gives \eqref{eq:family-selmer}
on this field at \(t\) with \(B=\sqrt b\), by nonnegativity of
the summands. It follows
also that \(4tr+1<q,d\) at sufficiently large \(x\).

At large \(x\) the rank interval for \(m\) gives
\(x^{1/2-\varepsilon}\le m\le x^{1/2+\varepsilon/2}\),
and the lower bound for \(k\) gives \(xm\le k\).
The upper bound for \(k\), together with \(n\ge x^{2+8\varepsilon}\),
gives \(k\le n^2x^{1/2-\varepsilon}\), since \(2\le x^{5\varepsilon}\)
eventually. These are the numerical rank bounds in
Sections~\ref{sec:integral} and~\ref{sec:fourier}.
For the trace comparisons, \(m\ge4\) eventually and
\(m^4\le16x^{2-\varepsilon/4}\le f\le k\le r\le x^5\)
for large \(x\). Moreover,
\[
 m^5 f\log b\le128 x^{9/2+(139/16)\varepsilon}
              \le128 kx^{-(21/16)\varepsilon}.
\]
Thus every fixed multiple of \(m^5 f\log b\) is smaller than \(k\)
beyond a threshold depending only on that multiple and \(\varepsilon\).

\paragraph{Level and sampling conditions at the classical host.}
\teamanchor{proof.classical-ranked-host-supply}%
We install the level primes and the two samples on the
same ranked field using the prime estimates and the splitting
conditions for the chosen pair.

Write \(K_d=\mathbf Q(\sqrt{-d})\) and \(Q=Q_Sqd\). The level
conditions specify two distinct
primes \(g_0,g_1\) in \([(kb)^8,2(kb)^8]\) with
\[
 g_i=\pi_i\overline{\pi_i},\qquad
 \pi_i\in\mathcal O_{K_d},\qquad \pi_i\equiv1\bmod Q\mathcal O_{K_d}.
\]
They also specify distinct primes \(p,p'\le Q^{C_p}\) for a fixed
\(C_p\), with \(p,p'\equiv1\bmod\operatorname{lcm}(15,Q)\).
These conditions are the level data of
\releasecite{Development/AbelianVarietyTheory/Definitions/GlobalTwisting/ThinScalesIffHostBound.lean}{101}{the chosen host}.
For the host supplied to the point construction take the auxiliary
twisting prime \(q_0\) to be the
\releasecite{Development/AbelianVarietyTheory/Definitions/GlobalTwisting/AuxiliaryPrimeMemAvoidedPrimes.lean}{75}{least prime}
congruent to \(1\bmod \ell_0Q\), where \(\ell_0\) is the least prime
not dividing \(Q(p-1)(p'-1)\). The host has unit ideal \(\mathfrak c_K\).
A prime divisor of \(Q(p-1)(p'-1)+1\) bounds
\(\ell_0\) by \(Q^{1+2C_p}+1\). If \(m_0=\ell_0Q\), the
full-interval prime theorem below at \(m_0^A\), for sufficiently
large fixed \(A\ge12\), has relative error of modulus less than \(1/2\).
The full-interval version of \eqref{eq:family-prime-multiplier}
gives a positive main term in the class \(1\bmod m_0\). Hence
\(q_0\le m_0^A\le Q^{A(3+2C_p)}\). Take
\(C_8\ge\max\{J,4,C_p,A(3+2C_p)\}\), so \(p,p',q_0\le Q^{C_8}\)
and \(Q_S\le C_8n\). Install this \(q_0\) with the other level primes.
The congruences distinguish \(q_0\) from \(p,p'\), since
\(\ell_0\) divides neither \(p-1\) nor \(p'-1\).
Since \(g_i\equiv1\bmod Q\) by the norm identity, all five primes
split completely in the cyclotomic field of modulus \(Q\) and hence
in \(E\) and \(F\). The logarithms of \(p,p',q_0\) are
\(O(x^{\varepsilon/5}(\log x)^2)\), which is eventually smaller
than \(x^{\varepsilon/3}\). Thus these primes are below both
\(\exp(x^{\varepsilon/3})\) and \((kb)^8\). Also
\(\log g_i\le8x^\varepsilon+8\log k+\log 2\le9x^\varepsilon\)
for large \(x\).

For each \(g\) let
\begin{equation}
\label{eq:classical-level-bad-primes}
 B_r(g)=\{g\}\cup
  \{\ell\ \text{prime}:\ell^{4v_\ell(g^a-1)}>g^a
                               \text{ for some }1\le a\le r\}.
\end{equation}
The disjointness \(B_r(g_0)\cap B_r(g_1)=\varnothing\) is part of
the selected level conditions. Thus for each residue characteristic
\(\ell\), one of the two installed \(g_i\) differs from \(\ell\)
and satisfies \(4v_\ell(g_i^a-1)\log\ell\le a\log g_i\) for
every \(1\le a\le r\).


\paragraph{Construction of the thick sample.}
\teamanchor{proof.thick-sample-construction}%
Put \(L_w=10\lceil405\log x\rceil+2\lceil30/\varepsilon\rceil+22\).
Fix \(C_{T_2}=16\).
Following Behrend's digit-and-sphere construction for three-term
progressions \cite[pp.~331--332]{f2-behrend}, choose a larger digit
base to treat every pair of coefficients of size at most \(L_w\).
We construct positive integer labels \(T_2\) with
\(4L_w s<n\) for every \(s\in T_2\), and with
\[
 |T_2|\ge\frac{n}{L_w}\exp\!\left(-16\sqrt{\log n\,\log L_w}\right).
\]
Let \(u=\log(n/(4L_w))\), \(v=\log(8L_w)\),
\(D=\lfloor\sqrt{u/v}\rfloor\), and
\(M_0=\lfloor e^{u/D}/(4L_w)\rfloor\). Since
\(n\asymp x^{2+8\varepsilon}\) and \(L_w=O_\varepsilon(\log x)\),
we have \(D\ge3\) and \(M_0\ge e^{u/D}/(8L_w)\) for large \(x\).
Take a sphere of maximal cardinality in the integer box
\(\{0,\ldots,M_0-1\}^{D}\). There are at most \(D M_0^2\)
possible squared norms, so that sphere has at least
\[
 \frac{M_0^{D}}{D M_0^2}
 \ge\frac{n}{L_w}
     \exp\!\left(-\log 4-\log D-Dv-2u/D\right)
\]
points. Here \(Dv\le\sqrt{uv}\), \(2u/D\le4\sqrt{uv}\), and
\(\log 4,\log D\le\sqrt{uv}\) for large \(x\). The asserted
cardinality bound follows from \(u\le\log n\) and
\(v\le4\log L_w\). The lower bound tends to infinity, so the
sphere has positive radius.

Encode each point by its base-\(\beta\) digits, where
\(\beta=4L_wM_0\). The values are distinct positive integers less
than \(\beta^D\le n/(4L_w)\). A relation
\(j(s_2-s_1)=i(s_3-s_1)\) with \(0<|i|,|j|\le L_w\) and
\(i\ne j\) would give a zero sum of powers of \(\beta\) whose signed
coefficients have absolute values less than \(2L_wM_0<\beta\).
Reducing successively modulo \(\beta\) shows that every coefficient
is zero. The three digit vectors would lie on one line, which meets
their sphere in at most two points. This excludes the relation on
distinct labels. The range \(4L_ws<n\) also puts the absolute value
of the relation below \(n\), so the same exclusion holds modulo \(n\).
The range and positivity imply \(4L_w|T_2|\le n\).
\releasecite{Formalism/FrobeniusOrbits/FrobeniusOrbitsProofs/ThickExists.lean}{181}{The released thick-set existence theorem}

The logarithmic loss \(\log L_w+16\sqrt{\log n\,\log L_w}\) is
\(o(\log x)\). Thus \(N_2=|T_2|\ge nx^{-\varepsilon}\) beyond a
common threshold.

\paragraph{Installing the classical field and its level data.}
\teamanchor{proof.classical-host-installation}%
The chosen \(q,d\) are larger than \(4tr+1\). Write
\(Q=Q_Sqd\) and take the small-prime splitting bound to be
\(\lfloor Z^{\varepsilon/5}\rfloor\). The quadratic Artin sign at
\(q\) describes the action on \(\sqrt{-q}\in E\). The dyadic modulus
is exact in the following sense. The unit with components
\(1+Q_S/2\bmod Q_S\) and \(1\bmod qd\) is \(1\bmod Q/2\) and
acts nontrivially on \(E\), since for \(h\ge2\) its dyadic component
is neither \(1\) nor \(-1\bmod Q_S\). Here \(Q_S\le n^2\) and
\(o^{64}=1\le2^h\). The prime-pair conditions give \(q\ne3\) and
the splitting assertions at the small odd primes. The condition
on odd divisors of \(Q_S\) is vacuous.

Choose an embedding \(\varphi_0\) with
\(\varphi_0(\sqrt{-d})=i\sqrt d\). In the coordinates of
\eqref{eq:artin-coordinates}, the automorphisms fixing
\(K_d=\mathbf Q(\sqrt{-d})\) are precisely the subgroup
\(P=\ker c\) labelling its CM type. The inertia computation in
\eqref{eq:classical-host-inertia} shows that the different
\(\mathfrak D_{E/K_d}\) is supported over the primes \(2,q\).
Indeed the inertia subgroup at \(d\) intersects \(P\) trivially.
Each of \(2,q\) splits into a conjugate pair in \(K_d\). Choose
one member of each pair and let \(\mathfrak b\) be the product of
the prime-power factors of the different above the chosen members.
The different is stable under \(\operatorname{Gal}(E/\mathbf Q)\)
by equivariance of the relative trace pairing. Consequently
\(\mathfrak D_{E/K_d}=\mathfrak b\bar{\mathfrak b}\),
\(\mathfrak b\) and \(\bar{\mathfrak b}\) are coprime, and
\(\mathfrak b\) is \(P\)-stable.

Take \(g_0,g_1,p,p'\) from
\hyperlink{proof.level-primes-short-intervals}{the level-prime construction}
and \(q_0\) from the prescribed least-prime rule. All five are
\(1\bmod Q\). The unit ideal \(\mathfrak c_K=\mathcal O_{K_d}\)
satisfies the ideal conditions by its norm \(1\), and the level
primes are coprime to \(Q_S\). Install these data,
\(\varphi_0,\mathfrak b\) and the integer labels \(T_2\) on the same
\(E\) and frame, with \(L_w,C_{T_2}\) as above and initial translate
zero. These thick labels satisfy the signed relation exclusion,
the density bound and \(4L_w|T_2|\le n\).
\releasecite{Development/AutomorphicArithmetic/CoprimeRanks/SecondSetNoRelation.lean}{13}{Field conditions and choices}
\releasecite{Development/AutomorphicArithmetic/CoprimeRanks/SecondSetNoRelation.lean}{92}{Extending the field by those choices}.

\paragraph{Selection of the top sample on the same host.}
\teamanchor{proof.classical-top-sample}%
Fix \(0<\varepsilon<1/52\) and \(J\ge1\) as in the classical
windows, and put \(\delta_1=x^{-\varepsilon/8}\).
Write \(h_p\in H\) for the \(H\)-coordinate of the arithmetic
Frobenius of an unramified prime. Use the host's inert class,
which excludes \(2,q,d\) and consists of primes for which the
two quadratic signs are both \(-1\). Put
\[
 X'=\exp(x^{\varepsilon/3}),\qquad
 \mathcal P=\{p<X':p\text{ prime in the inert class and }
                  \operatorname{ord}_H(h_p)\ge x^{\varepsilon/4}\}.
\]
Here the orbit step at a tested prime is \(h_p\). For a set
\(T\subset H\), an orbit is tested in the ordered positions
\(0,\ldots,\nu_p-1\), where \(\nu_p=\operatorname{ord}_H(h_p)\)
and position \(i\) represents \(t_0+ih_p\). The origin \(t_0\)
is the member of the orbit having least numerical host label.
Let \(B_p\) be the tops belonging to orbits that fail the test
specified below.

For any inert-class prime put
\[
 L_1(p)=\left\lceil
        \frac{\log((kx^{26})^4)}{\log p}\right\rceil,\qquad
 L_2=2\sup_{p\ {\rm in\ the\ inert\ prime\ class}}L_1(p),
\]
where the supremum of an empty collection of lengths is \(0\).
For every such prime \(p>Z^{\varepsilon/5}\) by small-prime
splitting, so \(\log p> (\varepsilon/(5J))\log x\).
At \(x\ge16\) the rank bound gives
\(\log((kx^{26})^4)\le124\log x\). Hence \(L_2\le L_*\),
where \(L_*=\lceil 2(620J/\varepsilon+1)\rceil\).
Use the following parameters of the orbit test:
\begin{gather*}
 h_G=2+8\varepsilon,\qquad K_G=\lceil8(h_G+2)/\varepsilon\rceil,
 \qquad r_G=K_G+1,\qquad U=\lceil1/\delta_1\rceil,\\
 \ell=(6r_G+1)U,\qquad
 M_0=\lceil(h_G+3)10^{6r_G+1}\log x\rceil,\qquad
 G'=(M_0+5)\ell,\\
 N_1=\left\lceil
       \max\{6\mathrm e\, G'\delta_1,(h_G+3)\log_2 x\}\right\rceil+1 .
\end{gather*}
The first test is that every cyclic window of \(L_2\) positions
contains at most \(K_G\) tops. Form \(\lfloor\nu_p/\ell\rfloor\)
full tiles of length \(\ell\) in the order starting from \(t_0\).
Number the \(6r_G+1\) cells of length \(U\) in each tile from \(1\).
A tile is marked if its odd cells are empty and each even cell
has exactly one top, at an even orbit position. The second test
requires a marked tile among the \(M_0+1\) consecutive full
tiles starting from any full tile, with the tile indices taken
cyclically. The third test is that every cyclic window of \(G'\)
positions contains at most \(N_1\) tops.
\releasecite{Formalism/LocalToGlobal/GlobalTheory/IsResidueClassZero.lean}{40}{The lengths \(L_1,L_2\)}
\releasecite{Formalism/LocalToGlobal/GlobalTheory/IsResidueClassZero.lean}{267}{The orbit origin}
\releasecite{Formalism/SharpRestriction/GrowthBounds/UnramifiedQuadraticSetting.lean}{239}{Orbit parameters}
\releasecite{Formalism/SubspaceCounting/GoodOrbitIntervalTileCondition.lean}{10}{The three tests}.

Enlarge the threshold in \(x\), uniformly in the host, so that
\(\delta_1\le1/100\), \(L_*<U\), and \(3G'\le x^{\varepsilon/4}\).
These hold eventually since \(G'=O_\varepsilon(x^{\varepsilon/8}\log x)\).
Then \(2L_1(p)\le L_2<U\) at each tested prime. At a marked
tile the consecutive selected positions differ by at least \(U+1\)
and the two end cells are empty. Every cyclic \(G'\)-window
contains a full marked tile. In fact the next full-tile start
is less than \(2\ell\) away. The successive \(M_0\) steps between
full-tile starts have total length less than \((M_0+1)\ell\),
including the possible terminal remainder. Including the marked
tile itself stays within \(G'=(M_0+5)\ell\). On a good orbit
of length \(\nu_p\), counting all \(G'\)-windows gives
\[
 \frac{\nu_p}{G'}\le |T\cap(t_0+\langle h_p\rangle)|
    \le \frac{\nu_p N_1}{G'}\le 20\delta_1\nu_p .
\]
For the last inequality, \(G'\delta_1\ge(M_0+5)(6r_G+1)\ge1\).
Here \(M_0\ge(h_G+3)10^{6r_G+1}\log x\), so
\((h_G+3)\log_2 x\le G'\delta_1\) and
\(N_1\le 6\mathrm e\,G'\delta_1+2\le20G'\delta_1\).

Choose each member of the actual \(H\) independently with
probability \(\delta_1\). Fix \(p\in\mathcal P\) and \(t\in H\).
Conditioning on \(t\in T\) fixes one position in its orbit and
leaves the other positions independent with the same probability.
If the first test fails, some \(L_2\)-window contains at least
\(K_G\) selected positions other than the fixed position. By a
union bound this has probability at most
\[
 \nu_p\binom{L_2}{K_G}\delta_1^{K_G}
 \le 2\binom{L_*}{K_G}x^{-2}.
\]
Here \(\nu_p\le n\le2x^{h_G}\) and \(\varepsilon K_G/8\ge h_G+2\).

In an unaffected full tile, an odd cell is empty with probability
\((1-\delta_1)^U>1/4\). This follows from
\(\log(1-\delta_1)\ge-\delta_1/(1-\delta_1)\) and
\(\delta_1 U\le1+\delta_1\). Each even cell contains at least
\(\lfloor U/2\rfloor\) even positions. The probability of exactly
one top there at an even position is at least
\(\lfloor U/2\rfloor\delta_1(1-\delta_1)^{U-1}>1/10\),
using \(U\ge100\) and \(U\delta_1\ge1\). Thus a full tile
avoiding the fixed position is marked with probability at least
\(c_*=10^{-(6r_G+1)}\). Among \(M_0+1\) consecutive full
tiles at least \(M_0\) avoid that position, and their marking
events are independent. The condition \(3G'\le\nu_p\) guarantees
that these full tiles are distinct. A union over the starting
tiles bounds the probability of failure of the second test by
\(\nu_p(1-c_*)^{M_0}\le\nu_p x^{-(h_G+3)}\).

If the last test fails, some \(G'\)-window contains at least
\(a=N_1-1\) selected positions other than the fixed one.
For one window this has probability bounded by
\(\binom{G'}{a}\delta_1^a\le(\mathrm e\,G'\delta_1/a)^a\)
when \(a\le G'\), and zero when \(a>G'\). The parameter \(a\)
is at least \(6\mathrm e\,G'\delta_1\) and \((h_G+3)\log_2 x\),
so the bound is at most \(2^{-a}\le x^{-(h_G+3)}\).
Taking the union over starting positions gives another
\(\nu_p x^{-(h_G+3)}\). Together the three estimates show
\[
 \Pr(t\in B_p\mid t\in T)\le C_bx^{-2},\qquad
 C_b=2\left(\binom{L_*}{K_G}+2\right).
\]
This calculation uses the exact orbit test at every tested prime.
\releasecite{Development/AutomorphicArithmetic/ThickFrames/NotGoodLe.lean}{84}{The released conditional probability bound}
\releasecite{Development/AutomorphicArithmetic/ThickFrames/orbitTerms.lean}{97}{Its three numerical terms}.

Put \(W_{\mathcal P}=\sum_{p\in\mathcal P}(\log p)/p\).
Summing the probabilities gives
\[
 \mathbf E\left[\sum_{p\in\mathcal P}\frac{\log p}{p}|B_p|\right]
 \le C_bx^{-2}\delta_1 n W_{\mathcal P}.
\]
If \(W_{\mathcal P}>0\), Markov's inequality bounds the chance
that the sum exceeds four times the right side by \(1/4\).
If \(W_{\mathcal P}=0\), the sum is zero. The sample size has
mean \(\delta_1 n\) and variance \(n\delta_1(1-\delta_1)\).
For \(\delta_1 n\ge16\), the chance that it differs from its
mean by more than \(\delta_1 n/2\) is at most \(1/4\).
Choose a single outcome satisfying both estimates, so that
\[
 \frac{\delta_1 n}{2}\le |T|\le 2\delta_1 n,\qquad
 \sum_{p\in\mathcal P}\frac{\log p}{p}|B_p|
     \le 4C_bx^{-2}\delta_1 nW_{\mathcal P}.
\]
The sole \(H_0=H\) coset contains every top and has prescribed
share \(1\). The images of the two sampling sets as CM places
are disjoint because their \(q\)-signs differ. Both embeddings
of their labels are injective, including the thick labels after
adding its translate, since they are distinct integers in
\([0,n)\). Together with the thick-set bound this gives
\begin{equation}
\label{eq:classical-sampling-size}
 0<\tfrac12 n x^{-\varepsilon/8}\le N_T=|T|\le2n x^{-\varepsilon/8},
 \qquad n x^{-\varepsilon}\le N_2=|T_2|\le n .
\end{equation}
\releasecite{Development/AbelianVarietyTheory/Definitions/GlobalTwisting/ThinScalesIffHostBound.lean}{36}{Host bad-top bound}
\releasecite{Development/AutomorphicArithmetic/ShortNorms/hostOfFieldOf.lean}{215}{Updating \(T\) on the chosen host}.

The kept set for the conductor-one target consists of the trivial
character, which lies in \(X_+\). The exceptional host character
has primitive associate \(\eta_q\) of conductor \(q\). The characters
of \(X_+\cup X_-\) are trivial on \(H_0=H\), so the bounded
complement consists of \(\eta_q,\eta_d,\eta_q\eta_d\).
Their centered estimates and \(4\log x\) bounds on the logarithmic
derivatives apply at this \(x\). These facts, the exact coset share
and the Selmer bound at the chosen \(t\) give the arithmetic
requirement on this host used in the common-family construction.
The real primitive associates with a \(q\)- or \(d\)-factor have
conductor at most \(8q^i d^j\) for the corresponding exponents \(i,j\).
The prime-pair size bounds put these conductors within the heights
in item~\ref{item:twist-boxes}, whose zero-free boxes therefore apply
to the actual host characters. The logarithmic derivative of the
sign character has modulus at most \(4\log x\le x^\varepsilon\).
The tail bound \eqref{eq:host-character-tail} below applies on the
same \(E\) to all primitive associates of its characters.
\releasecite{Formalism/SharpRestriction/RamificationFrobenius/TotientGaloisClasses/ExplicitFormulaErrorOfHostConditions.lean}{144}{Fourier character conditions}
\releasecite{Development/AutomorphicArithmetic/CharacterBoxes/Supply.lean}{42}{The character and Selmer requirement}.

For the passage to the realized records, note that
\(\sum_{p<X'}\log p/p\le2\log X'\). To see this take
\(N=\lceil X'\rceil-1\) and use
\(\lfloor N/p\rfloor\ge N/(2p)\) for \(p\le N\), together with
\(\sum_{p\le N}\lfloor N/p\rfloor\log p\le\log(N!)\le N\log N\).
Thus the weighted sum is at most
\(16C_b|T|x^{-2+\varepsilon/3}\le|T|x^{-2+\varepsilon}\)
once \(x^{2\varepsilon/3}\ge16C_b\). The fitted records
use the same Artin coordinates, top labels and least origins.
The same sum therefore gives the bad-top estimate in
\hyperlink{proof.selection-realized-samples}{the realized sample bounds}.
\releasecite{Development/AbelianVarietyTheory/Proofs/GlobalInterpolation/BadCharacterBound.lean}{50}{Transport to the realized record}.

For \(M>1\), at each sufficiently large \(x\) in the intervals Alp\"oge
\releasecite{Development/AutomorphicArithmetic/CharacterBoxes/SupplyChi.lean}{8}{chooses ranks, level primes, an Artin map and sampling sets}
\(T\subset H\) and \(T_2\subset H_0\) for \(E\). The choices
give \(|T_2|\ge x^{-\varepsilon}n\) and the distribution of \(T\) among
the cosets of \(H_0\) used in Section~\ref{sec:fourier}. The
\releasecite{Development/AutomorphicArithmetic/CharacterBoxes/EventuallyRankLeRpowOfLtOneDivFiftytwo.lean}{5}{rank bound} gives
the actual twist cutoff \(t\le\lfloor x^{400}\rfloor\), so
\eqref{eq:family-selmer} holds at \(t\) with bound \(\sqrt b\).
Write \(Q_S\) for the modulus of the compositum of the dyadic and
target-character fields. Since the primitive character \(\psi\)
factors through \(Q_S\), its conductor \(M\) divides \(Q_S\).
Also \(Q_S\mid\operatorname{lcm}(2^{h+2},M)\), so every prime factor of
\(Q_S\) divides \(2M\).
For the classical host, the only nontrivial characters left after
averaging over \(H_0=H\) are \(\eta_q,\eta_d,\eta_q\eta_d\). Each
satisfies \(|G_\chi(x)|\le x^{1/2+\varepsilon/2}\) and
\(|L'(1,\chi)/L(1,\chi)|\le4\log x\) for the same \(q,d\) at \(x\).
In particular the centering for \(c\) satisfies
\(|\ell_c|\le4\log x\le x^\varepsilon\) for large \(x\).
Together with the discriminant bound this gives
\(\log\operatorname{rd}(E)+|\ell_c|\le2x^\varepsilon\).
For \(M>1\) the primitive characters discarded after averaging
\releasecite{Development/AutomorphicArithmetic/CharacterBoxes/BoundedTw.lean}{45}{induce the nontrivial quadratic twists above}.
For them we use \(|G_\chi(x)|\le x^{1/2+\varepsilon}\), by weakening
the prime-pair estimate, and
\(|L'(1,\chi)/L(1,\chi)|\le4\log x\).

The modulus satisfies
\[
  \log(Q_Sqd)\le 2x^{\varepsilon/5}(\log x)^2 .
\]
Pullback by the Artin map gives a Dirichlet character modulo \(Q_Sqd\)
for each character of \(\operatorname{Gal}(E/\mathbf Q)\).
A primitive real nonprincipal associate has order two, and its
kernel in the cyclotomic Galois group fixes a quadratic field. For a
quadratic field \(\mathbf Q(\sqrt D)\) with squarefree \(D\), the
primitive Dirichlet character has conductor \(|D|\) when
\(D\equiv1\bmod4\) and \(4|D|\) otherwise. Its dyadic conductor
therefore divides \(8\) \cite[Ch.~V, Example~3.11]{f1milne-cft}.
By the coprime product decomposition of the residue group, a real
primitive associate is \(\rho\eta_q^i\eta_d^j\), with real \(\rho\)
of conductor dividing \(8M\). If \(q\) or \(d\) divides its conductor,
then \((i,j)\ne(0,0)\). Item~\ref{item:twist-boxes} excludes real zeros
\(\beta>1-1/400\) of that associate by taking the ordinate zero.
The other real nonprincipal associates have
conductor dividing \(8M\). Nonvanishing at \(1\) for this fixed finite
set gives a common real zero-free interval ending at \(1\).
For the classical host the real associates with a \(q\) or \(d\)
factor are exactly the three auxiliaries and, when \(h\ge1\), their
three twists by \(\nu\) of conductor \(8\). The only other real
nonprincipal associate is \(\nu\), when available. Thus one
\(\delta>0\) works for all real nonprincipal associates of all
these hosts, including those nontrivial on \(H\).

\teamanchor{proof.full-interval-prime-progressions}%
For integers \(m\ge2\), \((a,m)=1\), and real \(u\ge m^{12}\), put
\(\mathcal E_m(u,a)=\chi_m(a)u^{\beta_m}/\beta_m\) in the exceptional
case above and \(\mathcal E_m(u,a)=0\) otherwise. The full-interval
multiplier \(1-\mathcal E_m(u,a)/u\) lies between \(0\) and \(2\).
Thorner and Zaman's full-interval prime number theorem gives
\[
 \sum_{\substack{p\le u\\p\equiv a\bmod m}}\log p
 = \frac{u-\mathcal E_m(u,a)}{\varphi(m)}
   +O\left(\frac{u}{\varphi(m)}\,\mathcal R(u,m)\right),
\]
where
\[
 \mathcal R(u,m)=\exp\left(-c\frac{\log u}{\log m}\right)
   +\exp\left(-c\frac{(\log u)^{3/5}}{(\log\log u)^{1/5}}\right).
\]
The constant \(c>0\) and the implied constant are absolute and
effective. The bound on the multiplier gives the displayed error
uniformly in the reduced class \(a\).
\cite[arXiv v2, Corollary~1.4, Remark~1.5 and (1.7)]{f1thorner-zaman}

Put \(Q=Q_Sqd\) and \(\Psi_\chi(u)=\sum_{n\le u}\Lambda(n)\chi(n)\).
For \(u\ge y\ge\exp(x^{\varepsilon/3})/3\) and large \(x\), we have
\(u\ge Q^{12}\) by the bound on \(\log Q\). Weight the progression
estimate by the character induced by \(\chi\) modulo \(Q\) and sum
over the \(\varphi(Q)\) reduced classes. Orthogonality gives main
term \(u\mathbf1_{\chi=1}\). Since \(\chi_Q\) is real, the exceptional
term sums to zero except when the induced character equals \(\chi_Q\);
in that case the term is \(-u^{\beta_Q}/\beta_Q\).
In particular this term is absent for a principal or complex
associate. For any character \(\xi\) modulo \(Q\)
with primitive associate \(\xi^*\), the two \(L\)-functions satisfy
\[
 L(s,\xi)=L(s,\xi^*)\prod_{p\mid Q}(1-\xi^*(p)p^{-s}).
\]
The product is nonzero for \(\Re s>0\) \cite[(25.15.4)]{dlmf}.
Thus in the surviving case \(\beta_Q\) is a real zero of \(L(s,\chi)\)
and \(1-\beta_Q\ge\delta>0\) for a fixed \(\delta\).
The error from summing the residue classes is
\(O(u\mathcal R(u,Q))\). For each \(p\le\sqrt{u}\) the proper powers
of \(p\) contribute at most \(\log u\), and the omitted prime terms
with \(p\mid Q\) contribute at most \(\log Q\). Since
\(\beta_Q>1/2\) when it exists, the exceptional term divided by \(u\)
has modulus at most \(2u^{-\delta}\). Writing \(v=\log u\) now gives
\[
 \Psi_\chi(u)/u=\mathbf1_{\chi=1}
    +O(e^{-cv/\log Q}+e^{-c\sqrt v})
\]
with \(c>0\) uniform for the fixed target and \(\varepsilon\).

Take \(b(n)=\Lambda(n)\chi(n)\), \(a=\delta_\chi=\mathbf1_{\chi=1}\),
and \(R_\chi(u)=\Psi_\chi(u)-\delta_\chi u\).
By \eqref{eq:common-character-series}, the centering constant
\(C_b(1)\) is \(-L'(1,\chi)/L(1,\chi)\) for nonprincipal \(\chi\).
At the trivial character it is \(-1-\gamma\) by the
\hyperlink{proof.centered-mellin-residues}{zeta centering}.
The uniform estimate above gives \(\int_1^\infty |R_\chi(u)|u^{-2}\,du<\infty\)
at each fixed host. Thus \eqref{eq:common-ramp-tail} gives the
centered sum in the window normalization:
\[
  G_\chi(y)=-y\int_y^\infty R_\chi(u)\frac{du}{u^2}.
\]
The first error integrates to \(O((\log Q)e^{-c\log y/\log Q})\), with
\(\log y/\log Q\gg x^{2\varepsilon/15}/(\log x)^2\). The second is
\(O(e^{-c'\sqrt{\log y}})\), with \(\sqrt{\log y}\gg x^{\varepsilon/6}\).
Thus
\begin{equation}
\label{eq:host-character-tail}
  |G_\chi(y)|\le y\exp(-c_6x^{\varepsilon/8})
  \qquad\bigl(y\ge \exp(x^{\varepsilon/3})/3\bigr)
\end{equation}
for all primitive associates of the host characters, including \(\chi=1\).
The constant \(c_6>0\) and the threshold in \(x\) are uniform for the
fixed target and \(\varepsilon\). For the target of conductor one, the
square-root estimates apply to \(\eta_q,\eta_d,\eta_q\eta_d\), and the
large-argument estimate also applies to characters that are
nontrivial on \(H\).

\teamanchor{candidate.family}%
\textbf{AI-proposed candidate (uniform higher fine Selmer bounds on
windows).}
Fix the target \(\psi\), \(0<\varepsilon<1/52\), and an admissible
dyadic profile \(h\). The release's
\releasecite{Development/AutomorphicArithmetic/CharacterBoxes/ExistsPageN.lean}{44}{Selmer window assertion}
chooses \(J\ge1\) and \(X_0\) before \(X\) and the primes.
For every \(X\ge X_0\) and every pair with the prescribed sizes and
congruences at \(Z=X^{1/J}\), it gives \([W,W^2]\subseteq[Z,X]\).
At every real \(y\in[W,W^2]\), the field \(E(h(y),q,d)\) satisfies
\eqref{eq:family-selmer} with \(t=\lfloor y^{401}\rfloor\) and
\(B=\exp(y^\varepsilon/2)\), uniformly over all finite prime sets
and all choices of decomposition data.
\rhrange{The preceding window applies to every admissible prime pair at
a fixed primitive target, and permits the \(401\) cutoff at every
point of \([W,W^2]\).}
{At \(\psi=1\), \eqref{eq:family-selmer} through the \(400\) cutoff
on arbitrarily distant intervals \([U,U^{3/2}]\) for one analytically
selected pair per interval. The pure dyadic terms at \(3,7\) are
bounded on these intervals, and the other terms use the
\hyperlink{proof.selmer-classical-residual}{pointwise source estimate}
there.}

Poitou--Tate duality and \(K\)-theory give finiteness at fixed field
and twist. To compare with growth in a fixed tower, let \(K_n\)
be the degree-\(\ell^n\) layers of a fixed
\(\mathbf Z_\ell\)-extension \(K_\infty/K\), with no real places
when \(\ell=2\), and fix \(r\ge1\).
Lim gives
\[
 v_\ell\!\left(|WK_{2r}^{\rm et}(K_n)[\ell^\infty]|\right)
      =\mu(Y_r)\ell^n+\lambda(Y_r)n+O(1).
\]
Here \(WK_{2r}^{\rm et}[\ell^\infty]\) is the localization kernel
in degree two with coefficients \(\mathbf Z_\ell(r+1)\), and
\(Y_r\) is the norm limit of these kernels in the fixed tower.
The implicit constant and the invariants refer to that tower and
fixed \(r\), whereas the host \(E\) and the twist vary in our
estimate. \cite[arXiv v2, Proposition~5.1, p.~20]{f1lim-wild}
At \(r\asymp y^{400}\) and \([E:\mathbf Q]\asymp y^{2+8\varepsilon}\),
even the factor \((2r[E:\mathbf Q])^5\) on the right of
\eqref{eq:family-even-k-group} exceeds \(y^\varepsilon\).
The window at these scales uses a separate aggregate estimate.
The abelian equivariant Tamagawa number formula relates leading
terms to integral cohomology and a regulator in positive rank.
Bounding the order of the finite cohomology from such an identity
by estimating its factors separately would require lower bounds
for the regulator on the same fields and twists
\cite[Corollary~1.2]{f1burns-flach}
\cite[\S5.4, (26)]{f1soule-torsion}.

For \(\psi=1\), the
\hyperlink{proof.selmer-classical-residual}{classical Selmer estimate}
gives the bound at cutoff \(\lfloor x^{400}\rfloor\) throughout the
chosen prime-pair intervals. It supplies the arithmetic input to the
admitted family and the local division at the scale selected from
each window in Section~\ref{sec:assembly}. The quadratic-twist
estimates and the large-argument estimate \eqref{eq:host-character-tail}
on the same fields feed Section~\ref{sec:fourier}.
Section~\ref{sec:assembly} joins these inputs with the level and period
comparisons.
